\chapter{{Diachrony} {of} {\textsc{qualquier} }{revisited}}\label{sec:5}

The Old Spanish forms of \textit{cualquiera} were rather \textit{qualquier, -e, -a} or \textit{qual quier, -e, -a} (see \citealt{cp2009}, \citealt{Pato2012}, \citealt{Mackenzie2019}, among others). I will use the notation \textsc{qualquier} in reference to both Old and Modern Spanish. By contrast, the notation \textit{qualquier,-e,-a} refers exclusively to Old Spanish, and \textit{cualquier,-a} only to Modern Spanish. In this chapter, I will look at the syntactic and semantic change of the synthetic form of \textsc{qualquier}, the form that has been already reanalysed as a compound word, which has its origin in a free relative clause (see \citealt{Palomo1934}, \citealt{Lombard1938}–39/1947–48, \citealt{Meier1950}, \citealt{Rivero1988}, \cite{Espinosa2006}, \citealt{cp2009}, among others), illustrated again schematically as follows:

\ea \label{ex:key:272} %ex 272
\ea{} \label{ex:key:272-1}  [\textsubscript{MainClause} ….V\textsubscript{main}... [\textsubscript{CP-Rel.} \textbf{\textit{qualquier} }\textbf{[\textsubscript{C°}} \textbf{\textit{que}}] [\textsubscript{TP} \textbf{V\textsubscript{[subj.]}}]]] or
\ex{} \label{ex:key:272-2}  [\textsubscript{MainClause} [\textsubscript{CP-Rel.} \textbf{\textit{qualquier} }[\textbf{\textsubscript{C°}} \textbf{\textit{que}}] [\textsubscript{TP} \textbf{V\textsubscript{[subj.]}}]]] ...V\textsubscript{main} ...]
\z
\z

\ea \label{ex:key:273}
\ea \label{ex:key:273-1}
\gll Sepan quantos esta carta viren como {yo [...]} \textit{quito} \textit{e} \textit{escuso} \textit{el} \textit{su} \textit{portero} \textit{de} \textit{la} \textit{eglesia}, \textit{qual} \textit{quier} \textit{que} \textit{sea}, por todo tiempo de todo pecho [...]\\
know how.many this letter see how I remove and excuse the their usher of the church whatever may that be for all time of all payment\\
    \glt ‘May all that see this letter know how I […] I exempt and discharge the usher of the church, \textit{whoever (one wants that) it might be}, from all feudal taxes [...]’

    (13th c., Alfonso X, \textit{Doc. Cast.-León})
\ex \label{ex:key:273-2}
\gll E tan grandes eran las bozes delos hombres \& del ganado que \textit{qualquier} \textit{que} \textit{lo} \textit{viesse} hauria ende gran dolor \& muy gran piedad…,\\
and so great were the voices of.the men \& of.the cattle that whoever who it saw would.have from.it great sorrow \& very great pity\\
    \glt ‘and the voices of the men and of the cattle were so great that \textit{whoever} \textit{(one} \textit{wants} \textit{that)} \textit{s/he} \textit{sees} \textit{it} would have great pain and very great pity…,’

    (13th c., Anonymous, \textit{Gran Conquista})
\z
\z

\REF{ex:key:272} and \REF{ex:key:273}, respectively, show two kinds of embedding of the relative clause in the 13th century. In \REF{ex:key:272-1}, the relative clause introduced by \textit{qualquier} modifies an element of the main clause (the DP \textit{el su portero de la iglesia}), whereas in \REF{ex:key:272-2} the relative clause itself is the subject of the main clause. As will be shown in \ref{sec:5.2}, these are the most frequent types in the 13th century. Other types will be discussed in \ref{sec:5.3}. The separated form \textit{qual quier} (the prior stage of the synthetic form \textsc{qualquier}) will not be the focus in this section and the statistical analysis in \sectref{sec:5.1} and \sectref{sec:5.2}. The separated form \textit{qual quier} and the interposition form \textit{qual NP quier} was mostly used in the 13th century and disappeared from that period on (see \citealt{cp2009}, \citealt{Pato2012}), which is why I do not include this form in this study.

I will show in this section that the frequency of the use of the synthetic forms represented by \textsc{qualquier} in sentence types other than relative clauses as illustrated in \REF{ex:key:274} rose over time. For instance, the frequency of \textsc{qualquier} in main declarative clauses with indicative verbal mood as exemplified in a (schematically analysed in \ref{ex:key:274}) rises significantly (for other instances of sentence types, see \sectref{sec:5.2}):

\ea \label{ex:key:274}
\ea{} \label{ex:key:274-1} [\textsubscript{Main Declarative Clause} ... Verb\textsubscript{[Indic.]} ... [\textsc{qualquier} NP]]\textsubscript{subj}    (postverbal subject) or
\ex{} \label{ex:key:274-2} [\textsubscript{Main Declarative Clause} ... [\textsc{qualquier}][\textsubscript{subj} Verb\textsubscript{[Indic.]}]  (preverbal subject)
\z
\z

\ea \label{ex:key:275}
\ea \label{ex:key:275-1} \gll Por ende bien \textit{se} {avise} desto \textit{qualquier} varón.\\
for this well \textsc{refl} advise of.this \textsc{qualquier} man\\
    \glt ‘For this reason every man needs to get informed about this.’

    (14{th} c., López de Ayala, \textit{Libro rimado})
\ex{} \label{ex:key:275-2} [...], esto \textit{qualquier} \textit{sabría},[...]

‘[...], anybody would know it, [...]’

    (14{th} c., López de Ayala, \textit{Libro rimado})
\z
\z

As will be shown in this section, the sentence type and the verbal mood as well as other syntactic factors correlate with differences in the interpretation of \textsc{qualquier} (+/– EVAL). I will test the frequency changes of the type described in \REF{ex:key:273} vs. \REF{ex:key:274} and other changes quantitatively. The results of these tests will show how \textsc{qualquier} developed further on in time after it changed from a relative clause [\textsubscript{Rel.Cl.}= \textsubscript{D} \textit{qual} \textsubscript{V} \textit{quier} ...]\footnote{Rel.Cl. stands for relative clause.}, as discussed in \sectref{sec:5.2.1}, into an indefinite determiner [\textsubscript{D} \textit{qualquier/cualquier}] as in \textit{cualquier} \textit{persona} ‘any person’ (see \sectref{sec:5.2.1} on this grammaticalisation process).

The ultimate goal of investigating the diachronic change of \textsc{qualquier} is to see how it developed different meanings across time, such as the evaluative meanings ‘ordinary’, ‘common’, and ‘unremarkable’ and whether this development of meanings is related to the (morpho)syntactic structure and change of \textsc{qualquier.}

This chapter is organised as follows. \sectref{sec:5.1} introduces into the methodology and data selection chosen. The changes will be first described quantitatively in \sectref{sec:5.2} and verified statistically. \sectref{sec:5.3} suggests a syntactic and semantic analysis of the changes described in \sectref{sec:5.2}. A summary will be presented in \sectref{sec:5.4}.

\section{{Methodology} {and} {data}} \label{sec:5.1}

This section introduces into the methodology of this chapter.

\subsection{{Data} {selection}} \label{sec:5.1.1}

In order to investigate the development of \textsc{qualquier}, I used the diachronic subcorpus, \textit{Genre/Historical} of Mark Davies’ Corpus del Español (CdE\_G/H): \url{https://www.corpusdelespanol.org/hist-gen/}. This subcorpus contains more than 20,000 Spanish texts from the 1200s to the 1900s.\footnote{{There is another historical corpus, CORDE. This corpus is not annotated in any way and thus does not allow one to search for syntactic configurations. It is only useful for searching for exact phrases. Its main advantage, though, is that CORDE offers the possibility to limit the search to specific countries and registers across periods and varieties. CORDE contains 240 million words from 950 to 1974 and does not contain any oral data, relying exclusively on written texts from books and the press. It does not provide the relative frequency, which must be calculated manually. For these reasons I decided not to use it.}}

Different research questions formulated in \chapref{sec:5} have been answered using different data sets of the diachronic corpus CdE\_G/H. The research question concerning the development of \textsc{qualquier} and the correlation between parameters such as time and linguistic features (e.g. verbal position and sentence type) requires manual tagging of every example. For this research question a set of 150 examples per century was created, except for the 14th century, for which only 142 examples could be gathered due to the low number of data in this century in the diachronic corpus CdE\_G/H, which contains only 216 examples of \textsc{qualquier} in total, and except for the 19th century, for which 149 examples were used for reasons that will be explained below.

I searched the corpus for the string \textit{*ualqu*}. The asterisks serve as a wild card character for any letter or sequence of letters. The search was thus sensitive to forms such as: \textit{cualquier, cualquiera, qualquier, qualquiera, qualquiere} (see \citealt{cp2009}, \citealt{Pato2012} for graphical variants). The corpus data for answering the aforementioned research question does not include the separated form \textit{qual quier} and thus also excludes cases of interposition such as qual NP quier (see \citealt{Pato2012} on interposition). As I already explained, I use the generic notation \textsc{qualquier} to refer to all these variants.

The diachronic corpus contains data from the 13th to the 20th century. The corpus represents the centuries as follows: 1200s (1200–1300), 1300s (1300–1400), 1400s (1400–1500), 1500s (1500-1600), 1600s (1600–1700), 1700s (1700–1800), 1800s (1800–1900), 1900s (1900–2000) (\figref{fig:key:5}).\footnote{{In \sectref{sec:5.1} and \sectref{sec:5.2} of this chapter, I use screenshots from the output of the CdE search. In these screenshots, frequencies are indicated either as per million words as absolute numbers or both values are displayed. The information about the relative frequency is especially necessary where I compared frequencies of the whole corpus CdE, because the subcorpora used to represent every century or time period are not equal in CdE and the relative frequency allows us to compare the data between periods (see \citealt{Kellert2018questions} on relative frequency measurements in the CdE in Old Spanish).}}

\begin{figure}
\includegraphics[width=.6\textwidth]{figures/kellert-img003.png}
\caption{\label{fig:key:5} Screenshot of the corpus search\_CdE\_G/H.}
\end{figure}

For each century, at least 500 examples (with the exception of the 14th century) were selected for classification. I annotated the centuries in my data with ordinal numbers (e.g. 13th c.). For the first three centuries (i.e. 13th, 14th, and 15th centuries), I excluded some texts whose edition cannot be considered philologically reliable (see \citealt{OctaviodeToledo2017} and \chapref{sec:1}). Moreover, I excluded cases of coordination and ellipsis (see examples of coordination and ellipsis in Old Spanish in \sectref{sec:5.2}) so as to be able to annotate verbal properties of the examples with \textsc{qualquier} (see below, \sectref{sec:5.1.2}). After the exclusion of such occurrences of \textsc{qualquier}, I chose at least 150 random examples for each century (142 examples for the 14th c., as I have explained earlier), which is a sufficient number for a quantitative analysis. Due to the exclusion of coordination and ellipsis in the 19th c. (=1800s.), my sample comprised 149 examples.\footnote{{In order to make sure that my corpus is big enough to run statistical tests, I contacted} {Alexander Silbersdorff, Chair of Statistics at the University of Göttingen and Thomas Weskott (Germanistik, University of Göttingen). I especially thank Thomas Weskott for his advice on how to implement the model of logistic regression on “my data”.}} I call the data set acquired with this methodology “my data”. The data as well as the information about the selected texts (represented in percentages) and the cases that have been filtered out (e.g. coordination and ellipsis) can be found in the respective Excel file on GitHub: \url{https://github.com/olga-kel/Historical-Linguistics}.

The choice of the sample size of 500 examples per century represents a balance between the linguistic analysis and reliability of the philological work. The examples have been first checked for philological accuracy and around 150 examples have been selected for linguistic analysis. 150 examples per century provide a sufficiently large sample size to capture linguistic patterns and variations.

In addition, I performed a supplementary corpus-wide analysis, where possible. I indicated the relative frequencies per mil of the occurrences in the entire corpus. This method includes texts that are not necessarily philologically reliable, but it provides an overview of existing patterns across the entire corpus. As the corpus is only tagged for part-of-speech (POS), not for dependency relations such as subject and object, some linguistic variables like subject or object function of \textsc{qualquier} could not be identified automatically in the whole corpus.

I use examples from the whole corpus CdE\_G/H in order to make suggestions about previous diachronic stages of the grammaticalisation before \textsc{qualquier} was reanalysed as a compound word (see \sectref{sec:5.3}).
The distribution of words and texts per century and the division of text genre in the 20th c. (= 1900s) is shown in \tabref{tab:words_per_century}.


\begin{table}
\begin{tabularx}{.8\textwidth}{Xrr}
\lsptoprule
\textbf{Century} & \textbf{Words} & \textbf{\# Texts} \\
\midrule 1200s & 7,079,164 & 71\\
1300s & 2,667,810 & 50\\
1400s & 8,747,963 & 160\\
1500s & 17,774,762 & 323\\
1600s & 13,355,483 & 498\\
1700s & 10,324,328 & 176\\
1800s & 20,822,142 & 392\\
1900s-News & 5,144,631 & 6810 news articles \\
1900s-Lit & 5,144,073 & 850 novels + short stories\\
1900s-Oral & 5,113,249 & 2040 interviews \\
1900s-Acad & 5,138,077 & 2931 articles\\
\midrule \textbf{Total} & \textbf{101,311,682} & \textbf{13,926} \\
\lspbottomrule
\end{tabularx}
\caption{Distribution of words per century in the Historical Subcorpus of the Corpus del Español}
\label{tab:words_per_century}
\end{table}

\newpage
\subsection{{Classification} {of} {the} {examples}} \label{sec:5.1.2}
\largerpage[2]
As a result of the method described above, my data contain 150 examples with \textsc{qualquier} per century (except the 14{th} century and the 19{th} century). I classified these examples using different syntactic and semantic variables that were chosen based on my research questions. These variables can be clustered according to different groups represented as A–F. The graphic variation represents group A; \textsc{qualquier} as the modifier of an adjacent noun vs. \textsc{qualquier} alone (i.e. as bare Q without noun modification, or QN/NQ as nominal modification) is Group B; the syntactic function of \textsc{qualquier} (subject, direct object, indirect object, predicative noun, adverb) represents Group C; the position of \textsc{qualquier} with respect to the verb of which \textsc{qualquier} is an argument or which \textsc{qualquier} modifies if it is an adverbial (QV/VQ) represents Group D; the mood, tense, aspect, and finiteness of the verb of which \textsc{qualquier} is an argument or which \textsc{qualquier} modifies if it is an adverbial (indicative/subjunctive, infinitive/gerund/participle, perfective aspect\footnote{{Perfective aspect was chosen to determine whether} {\textit{cualquiera}} {can be used under perfective verbs in diachrony (see \chapref{sec:2} for the role of perfective verbs in synchrony).}}) constitutes Group E; and the sentence type and sentence mood of the sentence in which \textsc{qualquier} appears represent Group F. The latter group contains the clause type (main clause, subordinated clause) and the sentence mood: declarative or \textsc{qualquier} as antecedent of a relative or complement clause\footnote{{In the descriptive part, I describe a clause introduced by} {\textit{que}} {pretheoretically as a relative or complement clause. The distinction depends strongly on the analysis of} {\textit{quier}} {in} {\textit{qualquier}}.} introduced by \textit{que} (\textsc{qualquier} \textit{(XP) que} or \textit{Q que}), which I analysed as (free) relatives (FRel) or other type.

Before I present the data, it is important to distinguish between \textsc{qualquier} that appears in a main clause and \textsc{qualquier} that appears in a relative clause or other subordinate clauses, because this difference influences the classification of \textsc{qualquier}. As my interest is to understand the relation of \textsc{qualquier} to the elements that \textsc{qualquier} structurally depends on, I only classified the relation between \textsc{qualquier} and the structurally dependent verb. When \textsc{qualquier} introduced a relative clause, as in \REF{ex:key:272}, I classified \textsc{qualquier} with respect to the variables described above and used in \tabref{tab:key:17}–\tabref{tab:key:19} below (e.g. its syntactic function or its position to the verb) \textit{inside the relative clause} rather than inside the matrix clause. This classification is a logical consequence of my analysis of relative clauses, that is, \textsc{qualquier} is analysed as part of a relative clause (i.e. as a specifier of a relative clause) and that as such it is structurally connected to the elements inside the relative clause, such as the verb in subjunctive (see for instance \citealt{Quer2000} on the importance of the subjunctive in relative clauses).

\largerpage
I now proceed to go through the overall results concerning the different variables and to illustrate them with examples from the corpus.

\tabref{tab:key:12} shows the absolute frequencies for the variables of group A to D from the total of around 150 examples per century:

\begin{table}[b]
% \begin{tabularx}{\textwidth}{lrrrrrrrrrrrrr}
% \lsptoprule
% %\hhline%%replace by cmidrule{~~------------}
%  &  & \multicolumn{2}{c}{ \textbf{A}} & \multicolumn{3}{c}{ \textbf{B}} & \multicolumn{5}{c}{ \textbf{C}} & \multicolumn{2}{c}{ \textbf{D}}\\
% %\hhline%%replace by cmidrule{~~------------}
%  &  & \multicolumn{2}{c}{ \textbf{Orthography}} & \multicolumn{3}{c}{ \textbf{+/-} \textbf{noun} \textbf{modifier}} & \multicolumn{5}{c}{ \textbf{Syntactic} \textbf{function}} & \multicolumn{2}{c}{ \textbf{VQ}}\\
%  \textbf{Century} & \textbf{Occ} & \textbf{\textit{qual-}} & \textbf{\textit{cual-}} & \textbf{Bare} \textbf{Q} & \textbf{QN} & \textbf{NQ} & \textbf{Subj} & \textbf{DO} & \textbf{IO} & { \textbf{Copula}}
%
%  \textbf{Predicate} & \textbf{Adv} & \textbf{QV} & \textbf{VQ}\\
%  13th & 150 & 149 & 1 & 102 & 43 & 5 & 94 & 15 & 25 & 5 & 11 & 134 & 16\\
%  14th & 142 & 139 & 3 & 62 & 78 & 2 & 70 & 15 & 34 & 16 & 7 & 118 & 24\\
%  15th & 150 & 141 & 9 & 33 & 116 & 1 & 27 & 26 & 65 & 11 & 21 & 102 & 48\\
%  16th & 150 & 13 & 137 & 37 & 113 & 0 & 35 & 26 & 62 & 13 & 14 & 98 & 52\\
%  17th & 150 & 12 & 138 & 38 & 111 & 1 & 32 & 42 & 52 & 12 & 12 & 83 & 67\\
%  18th & 149 & 18 & 131 & 41 & 106 & 2 & 29 & 32 & 66 & 11 & 11 & 91 & 58\\
%  19th & 150 & 0 & 150 & 55 & 81 & 14 & 24 & 37 & 65 & 14 & 10 & 71 & 79\\
%  20th & 150 & 0 & 150 & 23 & 123 & 4 & 25 & 38 & 54 & 7 & 26 & 54 & 96\\
% \lspbottomrule
% \end{tabularx}

\begin{tabularx}{\textwidth}{Xrrrrrrrrr}
\lsptoprule
 & \textbf{13th} & \textbf{14th} & \textbf{15th} & \textbf{16th} & \textbf{17th} & \textbf{18th} & \textbf{19th} & \textbf{20th} \\
\midrule
\tablevspace
\multicolumn{8}{l}{\textbf{Orthography (A)}}\\
\midrule
qual- & 149 & 139 & 141 & 13 & 12 & 18 & 0 & 0 \\
cual- & 1 & 3 & 9 & 137 & 138 & 131 & 150 & 150 \\
\tablevspace
\multicolumn{8}{l}{\textbf{± noun modifier (B)}}\\
\midrule
bare Q & 102 & 62 & 33 & 37 & 38 & 41 & 55 & 23 \\
QN & 43 & 78 & 116 & 113 & 111 & 106 & 81 & 123 \\
NQ & 5 & 2 & 1 & 0 & 1 & 2 & 14 & 4 \\
\tablevspace
\multicolumn{8}{l}{\textbf{Syntactic function (C)}}\\
\midrule
Subj & 94 & 70 & 27 & 35 & 32 & 29 & 24 & 25 \\
DO  & 15 & 15 & 26 & 26 & 42 & 32 & 37 & 38 \\
IO & 25 & 34 & 65 & 62 & 52 & 66 & 65 & 54 \\
Copula predicate & 5  & 16 & 11 & 13 & 12 & 11 & 14 & 7 \\
ADV & 11 & 7 & 21  & 14 & 12 & 11 & 10 & 26 \\
\tablevspace
\multicolumn{8}{l}{\textbf{VQ (D)}}\\
\midrule
QV  & 134 & 118 & 102 & 98 & 83 & 91 & 71 & 54 \\
VQ & 16 & 24 & 48 & 52 & 67 & 58 & 79 & 96\\
\midrule
\textbf{total} & 150 & 142 & 150 & 150 & 150 & 149 & 150 & 150 \\
\lspbottomrule
\end{tabularx}


\caption{\label{tab:key:12} CdE\_cualquiera\_150 occ. per century, my data from A to D \REF{ex:key:1}}
\end{table}

In \REF{ex:key:276}, I provide examples for the variables of groups from B to D:

\ea \label{ex:key:276}
\ea  Bare Q\\
   \gll Et mando a \textit{qualquier} que touiere estas salinas que les de cadanno estos ocho caffizes de sal bien [...]\\
   and order to \textsc{qualquier} that has these salterns that to.them gives each.year these eight measures of salt well\\
    \glt ‘And I order anyone who has these salterns, that he should give them every year eight \textit{cahices} [measuring unit] of salt […]’

    (13{th} c., Alfonso X, \textit{Doc. Cast. - Castilla la Nueva})
\ex QN\\
    \gll Si puede tomar armas \textit{qualquier} \textit{persona}, sea la quistión primera.\\
if can take weapons \textsc{qualquier} person be the question first\\
    \glt ‘[The issue of] whether any person can take weapons, should be the first question.’

    (15{th} c., Rodríguez del Padrón, \textit{Triunfo de las donas})
\newpage
\ex NQ\\
   \gll y no parecería muy duro que hubiera quien ayudase a llenarle de un \textit{modo} \textit{cualquiera}.\\
and \textsc{neg} would.seem very hard that would.be someone helped to fill.it of a way \textsc{cualquiera}\\
    \glt ‘and it would not seem difficult [to find] that there would be someone who would help to fill it [i.e. the remaining space of an area on a newspaper page] \textit{in any (whichever) way}.’

    (19{th} c., Arenal, \textit{Beneficencia})
\newpage
\ex S=Subject\\
    \gll somos {del todo} enfermos e con mucha amargura; \textit{qualquier} {tentaçión} nos vençe, aunque non sea muy dura.\\
are completely sick and with much bitterness \textsc{qualquier} temptation us overcomes even.if \textsc{neg} is very hard\\
    \glt ‘we are totally weak and with much bitterness: any temptation wins over us, even if it is not hard [to resist].’

    (14{th} c., López de Ayala, \textit{Libro rimado})
\ex DO\\
    \gll Si puedo \textit{cualquier} \textit{sustancia} doctamente definir, y accidentes describir en la mayor importancia, …\\
if can \textsc{cualquier} substance scholarly define and accidents describe in the highest importance\\ 
    \glt ‘if I can define any kind of substance with erudition and describe any contingent property in its highest importance, ...’

    (17{th} c., Villamediana\textit{, Poesía})
\ex IO\\
    \gll [...] E a \textit{qualquier} delos {sobre dichos} que mostrare alguna destas escusas bien lo deuen oyr desque el casamiento sea hecho.\\
 {} and to \textsc{qualquier} of.the aforementioned that show some of.these excuses well it should hear once the marriage is made\\
    \glt ‘and \textit{from anyone} of the aforementioned that will \textit{show} some of these excuses, they well need to hear it before the marriage is made...’

    (13{th} c., Alfonso X, \textit{Siete partidas})
\ex QV\\
    \gll Sin dubda Dios perdona el pecado llorado; enpero \textit{qualquier} \textit{se} \textit{tema}, si puede aver ganado tal graçia; […]\\
without doubt God forgives the sin cried however \textsc{qualquier} \textsc{refl} fears if can have obtain such grace\\
    \glt ‘Without any doubt, God forgives the regretted sin; but anyone who may be frightened may have obtained such grace…’\\

    (14{th} c., López de Ayala, \textit{Libro rimado})
\newpage
 \ex  VQ\\
\gll Mi mamá política hacía trampas sin escrúpulos para \textit{ganar} \textit{cualquier} juego en que participara.\\
my mom political did tricks without scruples to win \textsc{cualquier} game in that participated\\
    \glt ‘My mother-in-law cheated in an unscrupulous way to win any game she was part of.’ 


    (20{th} c., Pisabarro, \textit{Del agua nacieron los sedientos})
\z
\z

\tabref{tab:key:13} shows all numerical values for the variables from Group E.

\begin{table}
% \begin{tabularx}{\textwidth}{XXXXXXXXX}
% \lsptoprule
% %\hhline%%replace by cmidrule{~~-------}
%  &  & \multicolumn{7}{c}{ \textbf{E}}\\
% %\hhline%%replace by cmidrule{~~-------}
%  &  & \multicolumn{7}{c}{ \textbf{Verbal} \textbf{Aspect/Mood/Finiteness}}\\
% &  & \multicolumn{2}{c}{ \textbf{Mood}} & \multicolumn{3}{c}{ \textbf{Inf}} & \multicolumn{2}{c}{ \textbf{Tense} \textbf{and} \textbf{Aspect}}\\
%  \textbf{Century} & \textbf{Occ} & \textbf{Ind} & \textbf{Sbjv} & \textbf{Imperative} & \textbf{Infinitive} & \textbf{Gerund} & \textbf{Participle} & \textbf{Perfective}\\
%  13th & 150 & 26 & 120 & 3 & 1 & 0 & 0 & 2\\
%  14th & 142 & 61 & 74 & 0 & 5 & 1 & 1 & 1\\
%  15th & 150 & 60 & 66 & 0 & 16 & 1 & 7 & 2\\
%  16th & 150 & 89 & 45 & 1 & 12 & 3 & 0 & 5\\
%  17th & 150 & 99 & 33 & 2 & 12 & 3 & 1 & 3\\
%  18th & 149 & 64 & 49 & 1 & 29 & 3 & 3 & 2\\
%  19th & 150 & 74 & 51 & 2 & 18 & 1 & 4 & 3\\
%  20th & 150 & 102 & 17 & 0 & 27 & 0 & 4 & 13\\
% \lspbottomrule
% \end{tabularx}

\begin{tabularx}{\textwidth}{Xrrrrrrrr}
\lspbottomrule
 & 13th &  14th &  15th &  16th &  17th &  18th &  19th &  20th \\
\midrule
\tablevspace
\multicolumn{8}{l}{\textbf{Mood (E)}}\\
\midrule
Indicative & 26   &  61   &  60   &  89   &  99   &  64   &  74   &  102 \\
Subjunctive & 120  &  74   &  66   &  45   &  33   &  49   &  51   &  17 \\
\tablevspace
\multicolumn{8}{l}{\textbf{Inf}}\\
\midrule
Imperative  & 3    &  0    &  0    &  1    &  2    &  1    &  2    &  0 \\
Infintive & 1    &  5    &  16   &  12   &  12   &  29   &  18   &  27 \\
Gerund  & 0    &  1    &  1    &  3    &  3    &  3    &  1    &  0 \\
\tablevspace
\multicolumn{8}{l}{\textbf{Tense/Aspect}}\\
\midrule
Participle & 0    &  1    &  7    &  0    &  1    &  3    &  4    &  4 \\
Perfective & 2    &  1    &  2    &  5    &  3    &  2    &  3  &  13\\
\midrule
\textbf{Total} & 150  &  142  &  150  &  150  &  150  &  149  &  150  &  150 \\
\lspbottomrule
\end{tabularx}

\caption{\label{tab:key:13} CdE\_cualquiera\_150 occ. per century from my data (E)}
\end{table}

In \REF{ex:key:277}, I provide examples for the variables of Group E. Recall here again that I only considered the mood of the main verb, only if \textsc{qualquier} was embedded inside a main clause. Otherwise, I considered the mood of the verb inside the relative clause, as in \REF{ex:key:272}.

\ea \label{ex:key:277}
\ea \label{ex:key:277-1}  Mood: indicative

    \gll Los actores le restaron elegancia al juego y \textit{desdeñaron} \textit{cualquier} posibilidad de gol.\\
the actors to.him reduced elegance to.the game and disregarded \textsc{cualquier} possibility of goal\\
    \glt ‘The actors downplayed the elegance of the game and discounted any possibility for a goal’

    (20{th} c., \textit{CR:PrLibre:98May4})
    
\ex \label{ex:key:277-2} Mood: subjunctive

    \gll E tan grandes eran las bozes delos hombres \& del ganado que \textit{qualquier} que lo \textit{viesse} hauria ende gran dolor \& muy gran piedad.\\
and so great were the voices of.the men \& of.the cattle that \textsc{qualquier} who it saw would.have from.it great pain \& very great pity\\ 
    \glt ‘and the voices of the men and of the cattle were so great loud whoever (one wants that) would see it would have great pain and very great pity...,’

    (13{th} c., Anonymous, \textit{Gran Conquista})
    
\ex \label{ex:key:277-3} Imperative

    \gll Y Él respondióle: -\textit{Demanda} \textit{qualquier} cosa que todo lo otorgaré.\\
and he answered.him ask \textsc{qualquier} thing that everything it will.grant\\
    \glt ‘And He answered him: Ask anything and I will give it to you.’

    (13th c., Anonymous, \textit{Siete sabios})
    
\ex \label{ex:key:277-4} Infinitive

\gll ... y que hemos dicho poderse \textit{hallar} en \textit{cualquiera} de los tres estilos.\\
{} y that have said can find in \textsc{cualquiera} of the three styles\\
    \glt ‘... and of which we have said that it is possible to \textit{find} it in \textit{any} of three styles’

    (18th c., Luzán, \textit{La poética})
    
\ex \label{ex:key:277-5} Gerund

    \gll […] o en Yahaurpampa, donde hubo la victoria de los Chancas, \textit{siendo} {cualquiera} de aquellos dos puestos más a propósito que el de Cacha […]\\
{} or in Yahaurpampa where was the victory of the Chancas being \textsc{cualquiera} of those two sites more to appropriate than the of Cacha\\ 
    \glt ‘or in Yahaurpampa, where there was the victory of the Chancas, any of those two (building) sites being more appropriate than the one in Cacha…’

    (16th c., Garcilaso de la Vega, \textit{Comentarios reales})
    
\ex \label{ex:key:277-6} Participle

    \gll Y si de \textit{qualquiera} pasion \textit{enpedidos} se hallan no sentencian en nada fasta ver se libres.\\
and if of \textsc{qualquiera} passion hindered \textsc{refl} find \textsc{neg} judge in nothing until see \textsc{refl} free\\
\glt ‘And if they find themselves hindered by any passion, they do not judge until they are free (from this passion).’

    (15th c., Diego de San Pedro, \textit{Cárcel})
    
\ex \label{ex:key:277-7} Perfect

    \gll \textit{Han} \textit{triturado} \textit{cualquier} intento de libertad intelectual de forma totalmente despiadada.\\
have crushed \textsc{cualquier} attempt of freedom intellectual of form totally ruthless\\
    \glt ‘They have destroyed any attempt of intellectual liberty in a totally ruthless way.’

    (20th c., Mortimer, \textit{Entrevista (ABC)})
\z
\z

The results for group F are as given in \tabref{tab:key:14}.

\begin{table}
% \begin{tabularx}{\textwidth}{XXXXXXX}
% \lsptoprule
% %\hhline%%replace by cmidrule{~~-----}
%  &  & \multicolumn{5}{c}{ \textbf{F}}\\
% %\hhline%%replace by cmidrule{~~-----}
%  &  & \multicolumn{5}{c}{ \textbf{Sentence} \textbf{Type/Mood}}\\
%  \textbf{Century} & \textbf{Occ} & \textbf{Subord.} & \textbf{Main} & \textbf{Decl} & \textbf{Q} \textbf{\textit{que}} & \textbf{Other} \textbf{Types}\\
%  13th & 150 & 146 & 4 & 4 & 127 & 19\\
%  14th & 142 & 111 & 31 & 37 & 75 & 30\\
%  15th & 150 & 122 & 28 & 37 & 67 & 46\\
%  16th & 150 & 109 & 41 & 48 & 63 & 39\\
%  17th & 150 & 96 & 54 & 59 & 33 & 58\\
%  18th & 149 & 115 & 34 & 39 & 53 & 57\\
%  19th & 150 & 98 & 52 & 56 & 37 & 57\\
%  20th & 150 & 89 & 61 & 80 & 15 & 55\\
% \lspbottomrule
% \end{tabularx}


\begin{tabularx}{\textwidth}{lYYYYYYYY}
\lsptoprule
&13th &  14th &  15th &  16th &  17th &  18th &  19th &  20th \\
\midrule
\tablevspace
\multicolumn{8}{l}{\textbf{Sentence type/mood (F)}}\\
\midrule
Subord.&146  &  111  &  122  &  109  &  96   &  115  &  98   &  89   \\
Main&4    &  31   &  28   &  41   &  54   &  34   &  52   &  61   \\
Decl&4    &  37   &  37   &  48   &  59   &  39   &  56   &  80 \\
Q \textit{que}&127  &  75   &  67   &  63   &  33   &  53   &  37   &  15 \\
Other types&19   &  30   &  46   &  39   &  58   &  57   &  57   &  55\\
\midrule
\textbf{Total} &150  &  142  &  150  &  150  &  150  &  149  &  150  &  150 \\
\lspbottomrule
\end{tabularx}
\caption{\label{tab:key:14} CDE\_cualquiera\_150 occ. per century, my data}
\end{table}

In \REF{ex:key:278}, I provide examples for the variables of Group F:

\ea \label{ex:key:278}
\ea \label{ex:key:278-1} Subordinate (Subord.)
    
    \gll \& diz que pero que eran tan grandes principes \textit{que} \textit{qualquier} dellos manternie una huerste yl darie recabdo que en aquella batalla murieron daquella guisa.\\
\& says that but that were so great princes that \textsc{qualquier} of.them would.maintain a army and=to.him would.give provision that in that battle died in=that manner\\
    \glt ‘and it is said that because they were such important princes that \textit{any} of them could maintain an army, and he would inform (them) that in that battle they died in that manner.’

    (13th c., Afonso X, \textit{General Estoria IV})
    
\ex \label{ex:key:278-2} Main
    
    \gll El término pájaro se aplica a \textit{cualquier} ave con capacidad para volar y de pequeño tamaño.\\
the term bird \textsc{refl} applies to \textsc{cualquier} bird with ability to fly and of small size\\
    \glt ‘The term bird can be applied to \textit{any bird} that is able to fly and small in size.’

    (20th c., \textit{Enc: Ave})
    
\ex \label{ex:key:278-3} Declarative (Decl)

    \gll El término pájaro se aplica a \textit{cualquier} ave con capacidad para volar y de pequeño tamaño.\\
    the term bird \textsc{refl} applies to \textsc{cualquier} bird with ability to fly and of small size\\
    \glt  ‘The term bird can be applied to \textit{any bird} that is able to fly and small in size.’

    (20th c., \textit{Enc: Ave})
    
\ex \label{ex:key:278-4} Q que

    \gll assi como que pechen omezillo / o otra pena \textit{qualquier} \textit{que} viere el alcalde que sera guisado.\\
    just.as as that pay fine / or other punishment \textsc{qualquier} that sees the mayor that will.be appropiate\\
    \glt ‘just as they pay for homicide/ or (receive) some other punishment whatever the mayor sees as appropriate.’

    (14th c., Anonymous, \textit{Leyes de estilo})
    
\ex \label{ex:key:278-5} other types

    \gll … pero quería ser rogado, \textit{para} \textit{que} no se le imputase jamás por traición \textit{cualquier} siniestro acaecimiento, sino por desgracia\\
{} but wanted be asked for that \textsc{neg} \textsc{refl} to.him attributed ever for treason \textsc{cualquier} unfortunate event but for misfortune\\  
    \glt ‘but he wanted to be asked (to do it), so that any unfortunate occurrence would ever be attributed to him for treason, but for misfortune.’

    (18th c., Bacallar y Sanna, \textit{Comentarios})
\z
\z

%After this introduction into the data classification, I will present a short introduction of statistical tests I used.

%\usepackage{comment}

%\begin {comment}

%\subsection{{Statistical} {tests} {and} {basic} {notions}} \label{sec:5.1.3}

%The aim of the statistical tests in this study is to compare frequencies of different constructions with a single period or to compare frequencies of a single construction in different periods. As these frequencies are never numerically identical, statistical tests are required to determine whether the difference in frequency is statistically significant. One possibility for measuring frequency differences is to use the \textit{chi square (}χ²) test. This test compares a set of observed frequencies O with a set of expected frequencies E (see \citealt{Baayen2008}: 77, \citealt{Author2018}: 64 and references therein): χ² = sum of (O–E)\textsuperscript{2} / E. Expected frequencies (E) are frequencies that we would obtain if the frequency difference was not statistically significant. In order to calculate them, we multiply the sum of observed frequencies and divide them by total number.

  %Let us imagine the following linguistic scenario. In the first period (Period 1), Construction 1 occurs 7 times and Construction 2 is observed 12 times. Moreover, in the second period (Period 2), Construction 1is encountered 13 times and Construction 2 occurs 8 times (\tabref{abc}).

%\begin{table}
%\begin{tabularx}{\textwidth}{XXX}
%\lsptoprule
%\hhline%%replace by cmidrule{~--}
% & \textbf{Construction} \textbf{1} & \textbf{Construction} \textbf{2}\\
 %\textbf{Period} \textbf{1} & 7 & 12\\
 %\textbf{Period} \textbf{2} & 13 & 8\\
%\lspbottomrule
%\end{tabularx}
%{\caption{\label{tab:} ~}Hypothetical data for demonstration of χ².}
%\end{table}

%We arrive at the expected frequency by calculating expected numbers for each individual cell (i.e. the frequencies we would obtain if there were no association between two variables, that is, between Construction and Period). We do this by multiplying the row sum by the column sum and dividing by the total number. The expected frequencies are given in brackets (\tabref{abc}).


%\begin{table}
%\begin{tabularx}{\textwidth}{XXXX}
%\lsptoprule
%\hhline%%replace by cmidrule{~---}
% & \textbf{Construction} \textbf{1} & \textbf{Construction} \textbf{2} & \textbf{Total}\\
%\textbf{Period} \textbf{1} & { 7}

 %(9.5) & { 12}

 %(9.5) & 19\\
%\textbf{Period} \textbf{2} & { 13}

 %(10.5) & { 8}

 %(10.5) & 21\\
%\textbf{Total} & 20 & 20 & Grand total=40\\
%\lspbottomrule
%\end{tabularx}
%{\caption{\label{tab:} ~}Expected and observed frequencies of hypothetical data. Bracketed numbers indicate expected frequencies.}
%\end{table}

%Now we calculate the formula of χ² : (7-9.5)²/9.5 = 0.6579. We continue to do this for each cell and add all the numbers. The figure we obtain is 2.5062, because 0.6579 + 0.6579 + 0.5952 + 0.5952 = 2.5062. We now calculate the degrees of freedom: (Number of Rows–1) × (Number of Columns–1)(2–1) × (2–1) 1×1=1 (degrees of freedom).

  %If the difference between observed and expected frequencies is large enough, we can reject the null hypothesis H\textsubscript{0}, which states that the frequency difference is not relevant (the so-called \textit{Hypothesis of Independence}). The alternative hypothesis to hypothesis H\textsubscript{0} is H\textsubscript{1}, which states that the frequency difference is relevant (see \citealt{Author2018}: 64 and references therein):

  %H\textsubscript{0} : χ² = 0 H\textsubscript{1} : x\textsuperscript{2} > 0

  %The χ² value for 1 degree of freedom must be equal or greater than 3.841, if we assume that it is 5\% likely that H\textsubscript{0} is true (see \sectref{sec:5.2} for the application of such a test) (see \citealt{Baayen2008}: 77, \citealt{Author2018}: 64 and references therein). If the obtained \textit{chi square} (χ²) is larger than 3.841, it implies that it is unlikely that our observed two frequencies have occurred by chance. The value for our example is 2.5062, which is below the critical value of 3.841. In other words, the two constructions (Construction 1 and Construction 2) are similar with respect to their frequency across time periods.

  %Another important notion of statistics used in this book, is the notion of \textit{standard deviation} (SD). The standard deviation measures the variation between a set of values (such as a set of observed frequencies of different constructions or one construction across time). A low standard deviation indicates that the values tend to be close to the mean (also called the expected value) of the set, while a high standard deviation indicates that the values are spread out over a wider range.

  %Let us assume we have a normal distribution of data, also called a Gaussian distribution, which has a shape of a bell with a peak at the average value (mean). The standard deviation is a measure of the data that disperse around the mean or the peak. In the following representation, about 68\% of the data points fall within 1 standard deviation of the mean, about 95\% of the data points fall within 2 standard deviations of the mean and about 99\% of the data points fall within 3 standard deviations of the mean. Note that only 3 out of 1,000 points will fall outside the area of 3 standard deviations. The larger the number of the standard deviation, the wider the spread is of the data around the mean.

%\begin{figure}
%{\caption{\label{fig:} ~}Normal distribution in http://stevegallik.org/cellbiologyolm_statistics.html.} 
%\includegraphics[width=.6\textwidth]{figures/kellert-img004.pdf}
%\end{figure}

%I also used \textit{Linear regression}, which is a model that can be used to identify the relationship or correlation between different variables by fitting a line to the observed data that best fits the data.


  %\figref{fig:key:16} shows this for an example from outside linguistics: the relation between cholesterol and body mass index (BMI). It shows that cholesterol and BMI correlate: the higher the BMI, the higher the cholesterol. \figref{abc} shows the estimated regression line superimposed on data points represented as a scatter diagram.

%\begin{figure}
%{\caption{\label{fig:} ~}The linear model for the relationship of Cholesterol and body mass index (BMI) in https://sphweb.bumc.bu.edu/otlt/mph-modules/bs/bs704\_correlation-regression/bs704\_correlation-regression\_print.html, accessed \citealt{November2020}.}
%\end{figure}

%The linear regression can be used in linguistics in order to test whether there is a linear correlation between the frequency of some linguistic construction and time or period, such as to test whether the frequency of postnominal \textsc{qualquier} changes across time or periods (see \sectref{sec:5.2.1} for such a test).


\section{Quantitative description of the results in my sample}\label{sec:5.2}

We can observe major quantitative differences in time with variables in almost every group in my sample, especially if we look at the frequency differences between the 13th century (1200–1300) and the 20th century (1900–2000). The following table contains variables that show a decrease in time:

\begin{itemize}
\item The subject function of \textsc{qualquier} decreases.
\item The use of bare \textsc{qualquier} (bare Q) without nominal modification decreases.
\item The preverbal position (QV) decreases.
\item The use of \textsc{qualquier} in clauses with subjunctive (Sbjv) decreases.
\item The use of \textsc{qualquier} in subordinated clauses decreases.
\item The use of \textsc{qualquier} \textit{que} (Q \textit{que}) decreases.
\end{itemize}


\begin{table}
\fittable{
\begin{tabular}{lrrrrrrr}
\lsptoprule
 {Century} & {Occ.} & {Subject} & {bare} {Q} & {QV} & {Sbjv} & { {Subord.}}  {clause} & {Q} {\textit{que}}\\
 \midrule
 13th & 150 & 94 & 102 & 134 & 120 & 146 & 127\\
& 100\% & 62.67\% & 68.00\% & 89.33\% & 80.00\% & 97.33\% & 84.67\%\\
\tablevspace
 14th & 142 & 70 & 62 & 118 & 74 & 111 & 75\\
& 100\% & 49.30\% & 43.66\% & 83.10\% & 52.11\% & 78.17\% & 52.82\%\\
\tablevspace
 15th & 150 & 27 & 33 & 102 & 66 & 122 & 67\\
& 100\% & 18.00\% & 22.00\% & 68.00\% & 44.00\% & 81.33\% & 44.67\%\\
\tablevspace
 16th & 150 & 35 & 37 & 98 & 45 & 109 & 63\\
& 100\% & 23.33\% & 24.67\% & 65.33\% & 30.00\% & 72.67\% & 42.00\%\\
\tablevspace
 17th & 150 & 32 & 38 & 83 & 33 & 96 & 33\\
& 100\% & 21.33\% & 25.33\% & 55.33\% & 22.00\% & 64.00\% & 22.00\%\\
\tablevspace
 18th & 150 & 29 & 41 & 91 & 49 & 115 & 53\\
& 100\% & 19.46\% & 27.52\% & 61.07\% & 32.89\% & 77.18\% & 35.57\%\\
\tablevspace
 19th & 149 & 24 & 55 & 71 & 51 & 98 & 37\\
& 100\% & 16.00\% & 36.67\% & 47.33\% & 34.00\% & 65.33\% & 24.67\%\\
\tablevspace
 20th & 150 & 25 & 23 & 54 & 17 & 89 & 15\\
& 100\% & 16.67\% & 15.33\% & 36.00\% & 11.33\% & 59.33\% & 10.00\%\\
\lspbottomrule
\end{tabular}
}

\caption{\label{tab:key:17} Decrease of frequency in my data with respect to variables: subject, bare Q without nominal modification, preverbal, subjunctive, subordinate clause, Q \textit{que}}
\end{table}

\figref{fig:key:8} shows the decrease in the frequency of these variables across time visually.

\begin{figure}
\includegraphics[width=\textwidth]{figures/kellert_img_fig5-2.png}
\caption{\label{fig:key:8} Decrease of frequency in my data with respect to the frequency of the variables: subject, bare Q without modification, preverbal, subjunctive, subordinate clause, Q \textit{que}}
\end{figure}

Although the decrease is not linear as the frequency of some variables increases in the 15th and 17th centuries (e.g. QV decreases until 17th c., increases in the 18th c., and decreases again from the 19th c. onwards), we can see an overall decrease from the 13th to the 20th century. This overall decrease is the focus of my statistical and linguistic analysis (see \sectref{sec:5.2.1} and \ref{sec:5.3}). It should be noted at this point that the interruption of the decrease in the data in the 15th and 17th centuries coincides with the importance of the 15th and 17th centuries in the history of Spanish language, which is subdivided into the 13th, 15th, 17th, and 19th centuries (see \sectref{sec:5.2.3}). The changes in frequencies in these centuries in \figref{fig:key:8} may reflect the grammatical changes in these centuries. The separated form \textit{qual quier,}-\textit{e,}-\textit{a} disappears from the 16th century onwards (\citealt{Pato2012}, \sectref{sec:2.6}) and that the 15th century can be considered the peak century when \textit{qualquier,-e,-a} grammaticalises completely. The changes in frequency around this peak century might also reflect the grammaticalisation of \textit{qualquier,} -\textit{e,} -\textit{a}. I will return to this hypothesis below when I discuss the different variables.

Let us look at some prototypical sentence from the 13th century that contains variables that decrease in their frequency in subsequent centuries.

\ea \label{ex:key:279} %ex 279
  \gll E tan grandes eran las bozes delos hombres \& del ganado que \textit{qualquier} \textit{que} lo \textit{viesse} hauria ende gran dolor \& muy gran piedad.\\
and so great were the voice of.the men \& of.the cattle that \textsc{qualquier} who it saw would.have hereof great pain \& very great pity\\
  \glt ‘and the voices of the men and of the cattle were so great that whoever (one wants that) s/he sees it would have great pain and very great pity.’

  (13th c., Anonymous, \textit{Gran Conquista})
\z

The example in \REF{ex:key:279} as the following structure. The relative clause \textit{qualquier que} \textit{lo viesse} is the subject of the main verbal phrase \textit{hauria ende gran dolor \& muy gran piedad} and, inside the relative clause headed by the complementiser \textit{que}, \textit{qualquier} is the subject of the verb \textit{viesse}. \textit{Qualquier} does not modify any overt noun, and it is used as a pronoun that introduces free relative clauses similar to \textit{quien} ‘who’ or \textit{el que} ‘(the one) who’; thus, the word order is QV. The verb \textit{viesse} inside the relative clause is used in the subjunctive. The subjunctive has a non-veridical interpretation, that is, one that lacks the presupposition that the proposition of the sentence with subjunctive is true (see \citealt{Quer1998,Quer2000}, \citealt{Becker2014}, and \chapref{sec:2}). The subjunctive does not allow a reference to a certain event, but it includes all possible situations in which someone has seen it.

The following example is similar, with \textit{qualquier} being used as a subject of the complement clause \textit{que le demandaua}. The main difference between \REF{ex:key:279} and \REF{ex:key:280} is that the verb of the complement clause is not in the subjunctive but in the indicative. The verb aspect of the complement clause is a durative aspect represented by imperfective form \textit{demandaua}, which has a generic interpretation in this example, due not only to the imperfective but also to the generic adverb \textit{siempre} ‘always’ (see \citealt{Quer2000}, and \chapref{sec:2} on the generic interpretation of imperfective aspect):

\ea \label{ex:key:280}
  \gll siempre daua a \textit{cualquier} que le demandaua.\\
  always gave to \textsc{cualquier} that him asked\\
  \glt ‘He always gave it to whoever asked him. = Every time someone asked him, he gave it to him.’

  (15th c., Fernando del Pulgar, \textit{Claros varones})
\z

The imperfective verb expressed in the embedded sentence in \REF{ex:key:280} has a function similar to that in \REF{ex:key:279}; that is, there is no reference to a particular event of someone giving something to someone. Instead, the imperfective verb together with the generic adverb \textit{siempre} ‘always’ expresses all possible events in the past of someone giving something to someone else when asked (see \citealt{Becker2014}).

I will argue in \sectref{sec:5.3} that the decrease of the subjunctive inside the relative clause is related to the degree of grammaticalisation of \textsc{qualquier}. The more transparent \textsc{qualquier} is as a compound containing a wh-word \textit{qual} and the verb \textit{querer}, the more cases we find with \textsc{qualquier} que + V\textsubscript{[subj]}. The less transparent \textit{quier} is as a verb, the fewer cases we find of the type \textit{cualquiera que} and for that matter, the more cases we find without the subjunctive. The abrupt fall of the subjunctive in the 14th century coincides with the decrease of the separated form \textit{qual quier} in the 14th century and its disappearance in the 15th.

Let us now discuss other changes. \textsc{qualquier} in relative clauses such as \textsc{qualquier} \textit{que} can only be realised in a preverbal position with respect to the embedded verb [Main clause ... (V\textsubscript{main}) ... [\textsc{qualquier} \textit{que} ... V\textsubscript{embed}] ... (V\textsubscript{main}) ...]. This follows from the analysis of \textsc{qualquier} as a specifier of a relative clause (see, e.g., \REF{ex:key:272} in this section).

Another prediction of the Grammaticalisation Hypothesis is that once the string \textit{qual quier} is no longer analysed compositionally as a wh-word and a verb, we should observe an increase of \textsc{qualquier} in sentence types other than relative clauses and the use of \textsc{qualquier} with a verbal mood other than the subjunctive. This is indeed what can be observed in the data in \tabref{tab:key:18}, which shows that all variables (i.e. postverbal position (VQ), indicative mood and infinitives, perfective aspect, main clause, declaratives, other sentence types, QN) tend to increase in time, especially if we compare the 13th century with the 20th.


\begin{table}
\fittable{
\begin{tabular}{lrrrrrrrrr}
\lsptoprule
 {Century} & {Occ.} & {VQ} & {Ind} & {Inf} & {Perf} & {Main} & {Decl} & {Other} & {QN}\\
 \midrule
 13th & {150} & {16} & {26} & {1} & {2} & {4} & {4} & {19} & {43}\\
& ~ & 10.7\% & 17.3\% & 0.7\% & 1.3\% & 2.7\% & 2.7\% & 12.7\% & 28.7\%\\
\tablevspace
 14th & {142} & {24} & {61} & {5} & {1} & {31} & {37} & {30} & {78}\\
& ~ & 16.9\% & 43.0\% & 3.5\% & 0.7\% & 21.8\% & 26.1\% & 21.1\% & 54.9\%\\
\tablevspace
 15th & {150} & {48} & {60} & {16} & {2} & {28} & {37} & {46} & {116}\\
& ~ & 32.0\% & 40.0\% & 10.7\% & 1.3\% & 18.7\% & 24.7\% & 30.7\% & 77.3\%\\
\tablevspace
 16th & {150} & {52} & {89} & {12} & {5} & {41} & {48} & {39} & {113}\\
& ~ & 34.7\% & 59.3\% & 8.0\% & 3.3\% & 27.3\% & 32.0\% & 26.0\% & 75.3\%\\
\tablevspace
 17th & {150} & {67} & {99} & {12} & {3} & {54} & {59} & {58} & {111}\\
& ~ & 44.7\% & 66.0\% & 8.0\% & 2.0\% & 36.0\% & 39.3\% & 38.7\% & 74.0\%\\
\tablevspace
 18th & {150} & {58} & {64} & {29} & {2} & {34} & {39} & {57} & {106}\\
& ~ & 38.7\% & 42.7\% & 19.3\% & 1.3\% & 22.7\% & 26.0\% & 38.0\% & 70.7\%\\
\tablevspace
 19th & {149} & {79} & {74} & {18} & {3} & {52} & {56} & {57} & {81}\\
& ~ & 53.0\% & 49.7\% & 12.1\% & 2.0\% & 34.9\% & 37.6\% & 38.3\% & 54.4\%\\
\tablevspace
 20th & {150} & {96} & {102} & {27} & {13} & {61} & {80} & {55} & {123}\\
& ~ & 64.0\% & 68.0\% & 18.0\% & 8.7\% & 40.7\% & 53.3\% & 36.7\% & 82.0\%\\
\lspbottomrule
\end{tabular}
}
\caption{\label{tab:key:18} Increase of frequency in my data with respect to variables: VQ, ind, inf, perf, main, decl, other sentence types, QN.}
\end{table}

\tabref{tab:key:18} shows the increase in the frequency of these variables across time visually (\figref{fig:key:9}).


\begin{figure}
\includegraphics[width=\textwidth]{figures/kellert_img_fig5-3.png}
\caption{\label{fig:key:9} ~Increase of frequency in my data with respect to variables: VQ, ind, inf, perf, main, decl, other sentence types, QN.}
\end{figure}

\largerpage
As with decreasing frequency (see \figref{fig:key:9}), the 15th and 17th centuries also stand out with respect to the increasing frequency. The frequency of QN increases strongly until the 15th century. The strong increase of QN until the 15th century can be correlated with the observation that the separated form \textit{qual quier,-e,-a} disappears after this century (see \sectref{sec:5.2.1}), which can be interpreted as an indication of completed grammaticalisation. The grammaticalisation of \textit{qual quier,-e,-a} until the 15th century correlates with the strong increase of the determiner \textit{qualquier} from the 13th to the 15th centuries.

Other variables (e.g. main clause, VQ, and declarative sentence type) increase until the 17th century, with an interruption in the 18th. The grammaticalisation of \textit{qualquier} correlates with the increase of different sentence types as an indefinite form no longer is a relative clause and thus no longer is restricted to a single sentence type.

\largerpage
Two prototypical examples that illustrate the increased frequency values of these variables are the following sentences with an imperative clause and a declarative clause. In both clauses, \textsc{qualquier} is followed by a noun (\textit{cosa} and \textit{varón}, respectively, with which it forms a single constituent). It is used in postverbal position with indicative verbal mood in a main clause and not introducing a relative or complement clause from which \textsc{qualquier} would be structurally dependent. Note that \textit{qualquier cosa} in \REF{ex:key:281-1}. is not an argument of the verb \textit{otorgaré} in the clause introduced by \textit{que}, but it is the argument of the verb in the main clause:

\ea \label{ex:key:281}
\ea \label{ex:key:281-1} \gll Y él respondióle: \textit{Demanda} \textit{qualquier} cosa que todo lo otorgaré.\\
and he answered.him ask \textsc{qualquier} thing that everything it I.will.grant\\
\glt ‘Ask anything and I will give it to you.’\footnote{This \textit{que} corresponds to what \citet[99]{koch1990gesprochene} call the ‘polyvalent \textit{que}’, which does not have a clearly subordinating function.}

    (13th c., Anonymous, \textit{Siete sabios})
    
\ex \label{ex:key:281-2} 

\gll {Por ende} bien \textit{se} \textit{avise} desto \textit{qualquier} varón.\\
therefore well \textsc{refl} take.heed of.this \textsc{qualquier} man\\
    \glt ‘Therefore any man should get well informed about this.’

    (14th c., López de Ayala, \textit{Libro rimado})
\z
\z

If the frequency change of increasing variables across time in \figref{fig:key:9} is statistically significant, something that needs still to be shown in the next section, we can conclude that the use of \textit{cualquiera} in relative clauses became less frequent, whereas its use in main clauses became more frequent over time:

\ea{} \label{ex:key:282}
  [MainClause ….V... [\textsc{qualquier} \textit{que} V\textsubscript{[subj.]}]] (most frequent in Old Sp., 13th c.)

  $\rightarrow$[Main Clause Verb\textsubscript{[Indic.]} \textsc{qualquier} NP] (most frequent in 20th-century Sp.)
\z

Let us now look at the variables more closely and test the change of frequency statistically.

\subsection{{Statistical} {analysis} {of} {relevant} {changes} {in} {my} {sample}} \label{sec:5.2.1}

This section introduces into the statistical analysis of relevant changes in my sample.

\subsubsection{{The decrease of QV order and Q que}} \label{sec:5.2.1.1}

\tabref{tab:key:19} shows the frequencies associated with different syntactic positions of \textit{cualquiera.} If we look at the overall frequency, it shows a continuous decrease of the preverbal position and a continuous increase of the postverbal position of \textsc{qualquier} in every century. Looking at the frequencies in every single century, we can observe a bigger decrease between the 14th and 15th centuries. This decrease might be motivated with the overall grammatical change of early Old Spanish into late Old Spanish which is considered to be close to the Modern Spanish (see \citealt{Ariza1998a}, \citealt{Penny1993}, among others). In particular, Old Spanish allowed several constructions of object preposing, which are, however, still attested in the 15th century, as in \REF{ex:key:285}, and even in the 17th century, as in \REF{ex:key:287}. By contrast, postverbal subjects have always been possible in Spanish, as in \REF{ex:key:283}. Another change from early Old Spanish to late Old Spanish (Modern Spanish) is the change in the determiner system which is directly related to some of the changes we observe with \textsc{qualquier}. The determiner system changed in this period of change from late Old Spanish to early Old Spanish/Modern Spanish (\cite{Mackenzie2019, Bartra, ranson2005variation, company2009sintaxis}). Latinised forms became less frequent in this period (\cite{Mackenzie2019, ranson2005variation}), and bare nouns decrease (\cite{Bartra, company2009sintaxis}). As already mentioned, the increase of determiners correlates with the increase of postnominal \textsc{qualquier} with articles, which contributes to the increasing use of postnominal \textsc{qualquier} as a modifier similar to adjectives.


\begin{table}
\begin{tabularx}{.8\textwidth}{Xr Yr Yr}
\lsptoprule
%\hhline%%replace by cmidrule{~~----}
 \textbf{Century} & \textbf{Total} & \multicolumn{2}{c}{\qquad \textbf{QV}} & \multicolumn{2}{c}{\qquad \textbf{VQ}}\\
%  \cmidrule(r){3-4}\cmidrule(l){5-6}
%   & \textbf{Occ.} & \textbf{Nr.} \textbf{of} \textbf{occ.} & \textbf{Rel.} \textbf{frequency} & \textbf{Nr.} \textbf{of} \textbf{occ.} & \textbf{Rel.} \textbf{frequency}\\
 \midrule
 13th & 150 & 134 & 89.33\% & 16 & 10.67\%\\
 14th & 142 & 118 & 83.10\% & 24 & 16.90\%\\
 15th & 150 & 102 & 68.00\% & 48 & 32.00\%\\
 16th & 150 & 98 & 65.33\% & 52 & 34.67\%\\
 17th & 150 & 83 & 55.33\% & 67 & 44.67\%\\
 18th & 149 & 91 & 61.07\% & 58 & 38.93\%\\
 19th & 150 & 71 & 47.33\% & 79 & 52.67\%\\
 20th & 150 & 54 & 36.00\% & 96 & 64.00\%\\
\lspbottomrule
\end{tabularx}
\caption{\label{tab:key:19} Distribution of \textit{cualquiera} NP in preverbal (QV) or postverbal (VQ) position in my data}
\end{table}

The following examples demonstrate minimal pairs (QV and VQ) in my data:

\ea \label{ex:key:283}
  13th century
\ea \label{ex:key:283-1} 
QV 

\gll Et \textit{qualquiere} que les contra esto \textit{fuesse} a el me tornaria por ello.\\
and \textsc{qualquiere} that to.them against this would.go to him myself would.turn for it\\
\glt ‘And I would go against whoever would be against it.’

      (Alfonso X, \textit{Doc. Cast. – Murcia})
      
\ex \label{ex:key:183-2} 
VQ

\gll \& por la altura del polo quando \textit{fuere} sabudo \textit{qualquier} {dellos}.\\
 \& by the height of.the pole when will.be known \textsc{qualquier} of.them\\
\glt ‘and for the height of the pole, as soon as any of them is known’

      (Alfonso X, \textit{Albateni})
\z
\z

\ea \label{ex:key:284}
  14th century
\ea \label{ex:key:284-1} 
QV

\gll ca \textit{qualquier} que lo faga, çierto sea de ser {mal andante} ante que muera.\\
 for \textsc{qualquier} that it does certain will.be of being unfortunate before he dies\\
\glt ‘because anyone that would do it, would be certainly unfortunate before he dies.’

(Anonymous, \textit{Libro del Caballero Zifar})
      
\ex \label{ex:key:284-2} 
VQ

\gll  \& el señor {puede} {a} \textit{qualquier} destos mayordomos ante que se despidan del prender los \& tener los presos:\\
 \& the lord can to \textsc{qualquier} of-these stewards before that \textsc{refl} resign from.the seize the \& hold the prisoners\\
\glt ‘and the lord can arrest and imprison any of these stewards before they resign from their services to him:’

      (Anonymous, \textit{Leyes del estilo})
\z
\z

\ea \label{ex:key:285}
  15th century
\ea \label{ex:key:285-1} QV  

\gll Y si \textit{de} \textit{qualquiera} \textit{pasion} enpedidos \textit{se} \textit{hallan} no sentencian en nada falta ver se libres.\\
and if of \textsc{qualquiera} passion hindered \textsc{refl} find \textsc{neg} judge in nothing lacks see \textsc{refl} free\\ 
\glt ‘And if they find themselves hindered by any passion, they do not judge until they are free (from this passion).’

      (Diego de San Pedro, \textit{Cárcel})
      
\ex \label{ex:key:285-2} 
VQ

\gll  E todos los que en aquellos tiempos vinieron ale \textit{demandar} \textit{qualquier} {cargo}.\\
and all the who in those times came to=him demand \textsc{qualquier} {cargo}\\
\glt ‘and all those who, at that time, came to ask him from any (kind of) office’

      (Fernando del Pulgar, \textit{Claros varones})
\z
\z

\newpage
\ea \label{ex:key:286}
  16th century
\ea \label{ex:key:286-1} 
QV

\gll Consideraba \textit{cualquiera} ocupaçión que la \textit{podía} estorbar, […]\\
considered \textsc{cualquiera} occupation that it could hinder\\
\glt ‘I considered any activity that could bother her, ...’

(Villalón, \textit{El Crotalón})
      
\ex \label{ex:key:286-2} 
VQ 

\gll Un desconfiar de sí y un estar siempre temiendo que \textit{podrá} exceder al mío \textit{cualquiera} \textit{mérito} ajeno.\\
a distrust of oneself and a being always fearing that could exceed to=the mine \textsc{cualquiera} merit foreign\\
\glt ‘A lack of trust in oneself and a constant fear that any merit of someone else will exceed my own.’

(Juana Inés de la Cruz, \textit{Inundación Castálida})
\z
\z

\ea \label{ex:key:287}
  17th century
\ea \label{ex:key:287-1} 
QV 

\gll De don Álvaro, por fe, \textit{cualquier} cosa \textit{creeré}, en razón de la que vi.\\
of Mr. Álvaro, by faith \textsc{cualquier} cosa will.believe in reason of the that saw\\
\glt ‘From don Álvaro, in faith, I would \textit{believe} \textit{anything}, based on what I saw.’

(Castro, \textit{Malcasados})
      
\ex \label{ex:key:287-2} 
VQ 

\gll y no \textit{basta} \textit{cualquiera} para mandar\\
and \textsc{neg} suffices \textsc{cualquiera} for command\\
\glt ‘and not anyone is appropriate to be in charge of command’

(Mira de Amescua, \textit{Examinarse de rey})
\z
\z

\ea \label{ex:key:288}
  18th century
\ea \label{ex:key:288-1} 
QV 

\gll pues es sabido y consta a todo el mundo que \textit{qualquier} diarista \textit{es} un hombre infalible y no puede engañarse en materia alguna\\
for is known and is.evident to all the world that \textsc{qualquier} journalist is a man infallible and \textsc{neg} can be.mistaken in matter any\\
\glt ‘because it is known and everybody is aware of the fact that any journalist is an infallible man and cannot be fooled in any subject…’

(Forner, \textit{Gramáticos})
      
\ex \label{ex:key:288-2} 
VQ

\gll pues la semilla de que \textit{se} {forma} \textit{cualquiera} planta, necesariamente vino de otra planta de la especie misma\\
since the seed of that \textsc{refl} forms \textsc{cualquiera} plant necessarily came from another plant of the species same\\
\glt ‘because the seed from which any plant grows came necessarily from another plant of the same species’

(Feijoo, \textit{Teatro crítico})
\z
\z

\ea \label{ex:key:289}
  19th century
\ea \label{ex:key:289-1} 
QV 

\gll \textit{Cualquiera} {me} {creerá} tu hermanito menor.\\
\textsc{cualquiera} me will.believe your little.brother younger\\
\glt ‘Anybody will believe that I am your younger brother.’

(Tamayo y Baus, \textit{No hay mal})
      
\ex \label{ex:key:289-2} 
VQ 

\gll Luego, salía, \textit{y} \textit{dirigíase} lentamente por las calles hacia \textit{cualquier} \textit{café}, en donde les escribía a su madre y a la novia.\\
then went.out and was.heading slowly through the streets towards \textsc{cualquier} café, in where to.them wrote to his mother and to the girlfriend\\
\glt ‘Then he left, and headed slowly through the streets towards any café, where he would write to his mother and his fiancée.’

      (Trigo, \textit{En la carrera})
\z
\z

\ea \label{ex:key:290}
  20th century
\ea \label{ex:key:290-1} 
QV 

\gll \textit{Cualquier} {persona} {puede} disfrutar una biblioteca de ensueño.\\
\textsc{cualquier} person can enjoy a library of dream\\
\glt ‘Anyone can enjoy a dream library.’

(Boo Juan Vicente, \textit{España:ABC: CONTEXT +})
      
\ex \label{ex:key:290-2} 
VQ 

\gll Esto lo \textit{puede} hacer \textit{cualquiera}.\\
this it can do \textsc{cualquiera}\\
\glt ‘Anybody can do this.’

(@ 1)
\z
\z

The overall decrease in the order QV is parallel to the overall decrease of the occurrence of \textsc{qualquier} \textit{que} as \tabref{tab:key:20} shows. Looking more closely into the frequency distribution of different centuries, the decrease is drastic from the 13th to the 15th century, which might again be correlated with the grammatical change of early and late Old Spanish and the grammaticalisation of \textit{qual quier} as discussed with respect to \figref{fig:key:9}.


\begin{table}
\fittable{
\begin{tabular}{lrrr}
\lsptoprule
 \textbf{Century} & \textbf{Analysed} \textbf{examples/occurences} \textbf{(occ.)} & \textbf{Nr.} \textbf{of} \textbf{occ.} & { \textbf{Rel.} }  \textbf{frequency}\\
 \midrule
 13th & 150 & 127 & 84.67\%\\
 14th & 142 & 75 & 52.82\%\\
 15th & 150 & 67 & 44.67\%\\
 16th & 150 & 63 & 42.00\%\\
 17th & 150 & 33 & 22.00\%\\
 18th & 149 & 53 & 35.57\%\\
 19th & 150 & 37 & 24.67\%\\
 20th & 150 & 15 & 10.00\%\\
\lspbottomrule
\end{tabular}
}
\caption{\label{tab:key:20} Distribution of \textsc{qualquier} \textit{que} in my data}
\end{table}

\figref{fig:key:10} visualises the change in the frequency of preverbal \textsc{qualquier} (QV) and \textsc{qualquier} que (\textsc{Q} \textsc{que}) in \tabref{tab:key:20}.

\begin{figure}
\includegraphics[width=\textwidth]{figures/kellert_fig_10.png}
\caption{\label{fig:key:10} Decrease of QV and Q \textit{que} (=\textsc{qualquier} \textit{que}) in my data.}
\end{figure}

The question is whether the overall change is statistically significant. On the null hypothesis, the differences in frequencies per century are not significant; the difference in observed values could be just a random effect of different samples. If the null hypothesis is correct, the frequencies per century should be the same as or very close to the mean. On the intuitive level, the null hypothesis is not correct due to large differences in the data. We know that the maximal frequency is 127 (13th c.) and the minimal frequency is 15 (20th c.). The mean of all frequencies per century is 59. However, the standard deviation (SD) of \textsc{qualquier} \textit{que} is 32, which is almost half as large as the mean, which is far too high to reflect a normal distribution of the data. The larger the number of the standard deviation, the wider the spread is of the data around the mean. The data in \tabref{tab:key:20} cannot be distributed normally. This is also shown by the fact that the data in \figref{fig:key:10} does not show a bell curve comparable to that of a normal distribution (\cite{Baayen2008}).

The same difference obtains for the QV/VQ Variable. The maximal frequency is 134 (13th c.) and the minimal frequency is 54 (20th c.). The mean of all frequencies per century is 94. However, the standard deviation of QV is 24, which is far too high to reflect a normal distribution of the data. Thus, even from a descriptive level, the difference is high enough to conclude that the time or century matters for the variables \textsc{qualquier} \textit{que} and QV and that these two variables are correlated. The descriptive data is summarised in \tabref{tab:key:21}.


\begin{table}
\begin{tabularx}{\textwidth}{XXXXXX}
\lsptoprule
\multicolumn{2}{c}{ Mean} & \multicolumn{2}{c}{ Variation (s²)} & \multicolumn{2}{c}{ Standard deviation (Std.)}\\
\cmidrule(lr){1-2}\cmidrule(lr){3-4}\cmidrule(lr){5-6}
 {QV} & {Q} {\textit{que}} & {QV} & {Q} {\textit{que}} & {QV} & {Q} {\textit{que}}\\
 \midrule
 93.875 & 58.75 & 564.375 & 1011.5 & 23.756578 & 31.8040878\\
\lspbottomrule
\end{tabularx}
\caption{\label{tab:key:21} Descriptive statistics of QV and of Q \textit{que}}
\end{table}

In order to test the null hypothesis, namely that time (here: century) has no influence on the frequencies of QV and on \textit{cualquiera} used in relative clauses (i.e. \textsc{qualquier} \textit{que}), I used a linear regression model as shown, in \figref{fig:key:11} and \figref{fig:key:12}. These figures show that time plays a role in the frequencies of both constructions; that is, the use of both QV and \textsc{qualquier} \textit{que} changes statistically significantly over time. The regression output below shows that the Century variable has a significant influence on the syntactic structure (QV) due to the t-value, which correlates to a significance value (t = –4.68, p < .01) in \figref{fig:key:11}. This shows that Century has an influence on the syntactic structure \textsc{qualquier} \textit{que} due to the high significance value (t = –3.940, p < .01).\footnote{{To express the result in statistical terminology, I computed a linear regression on the logit-transformed relative frequencies of QV and} \textsc{qualquier}{ \textit{que}} {to test the null hypothesis that the variable Century has no influence on the frequency of QV and/or} \textsc{qualquier} {\textit{que}}{. The linear regression showed that Century has a significant influence on QV (t = –4.68, p < .01) and} \textsc{qualquier} {\textit{que}} {(t = –3.940, p < .01).}} The output of the linear regression is given in \figref{fig:key:12}.


\begin{figure}
\begin{lstlisting}
Coefficients:
            Estimate Std. Error t value Pr(>|t|)
(Intercept)  2.5804      0.4449   5.800  0.00115 **
Century     -0.4122      0.0881  -4.678  0.00340 **
\end{lstlisting}

\caption{\label{fig:key:11}The output of the linear regression from R with respect to QV.}
\end{figure}


\begin{figure}
% \begin{tabularx}{\textwidth}{X}
% \lsptoprule
\begin{lstlisting}
Coefficients:
            Estimate Std. Error t value Pr(>|t|)
(Intercept)   2.0768     0.6924   2.999  0.02403 *
Century      -0.5402     0.1371  -3.940  0.00763 **
\end{lstlisting}
% \lspbottomrule
% \end{tabularx}
\caption{\label{fig:key:12} The output of the linear regression from R with respect to \textsc{qualquier} \textit{que.}}
\end{figure}

I also ran a \textit{chi square} test of independence to examine the relation between verbal position and century. The relation between these variables was significant, $\chi^2$ (degree of freedom = 1, 91.1854), p < 0.00001; that is, VQ is more likely to appear later than QV.

My results confirm the general observation about the decrease of \textsc{qualquier} \textit{que} in Old Spanish in the literature (\citealt{cp2009}) on the basis of statistical tests. Additionally, I have shown that it is not only \textsc{qualquier} \textit{que} that decreases but also QV; and for that matter, VQ increases. In \sectref{sec:5.3}, we will see what consequencies this change in frequency brings for the semantic interpretation of \textsc{qualquier} and how this change in frequency follows from grammaticalisation.

\subsubsection{{Frequency distribution of QN/NQ}} \label{sec:5.2.1.2}

My sample has low numbers of postnominal \textsc{qualquier} (henceforth NQ) (ranging between 0 and 4 occurrences out of 150 occurrences. in every century). Recall that my data do not contain coordination and ellipsis data. As NQ appears a lot in coordination and elliptical constructions, as will be shown in \sectref{sec:5.2}, NQ is not sufficiently represented in my data. In order to test the change of NQ and QN including elliptical sentences and coordination, I used the whole diachronic corpus.

As we can see just by the observation of the percentage numbers of NQ and QN per every century in the figure and table below, NQ and QN have more or less the same distribution in the 13th century: around 50\% (NQ having 52\% and QN 48\%). This frequency difference changes drastically in the 14th century, as shown in \tabref{tab:key:22}; that is, QN increases, whereas NQ decreases, except in the 18th and 19th centuries, where QN decreases and NQ increases. This difference in frequencies needs an explanation, which will be provided in \sectref{sec:5.3}.


\begin{figure}
\includegraphics[width=\textwidth]{figures/kellert_fig_13.png}
\caption{\label{fig:key:13} Distribution of percentages of QN/NQ per century in CdE}
\end{figure}

\begin{table}
% \begin{tabularx}{\textwidth}{XXXXXX}
%  \textbf{Century} &  \textbf{QUALQUIER N in \%} &   \textbf{QUALQUIER N in nr. of Frequency} &
%  \textbf{N QUALQUIER in \%} &  \textbf{N QUALQUIER in nr. of Frequency} &  \textbf{Total nr. of Frequency of QUALQUIER}\\
%  13th & 48 & 168 & 52 & 185 & 353\\
%  14th & 83 & 106 & 17 & 21 & 127\\
%  15th & 90 & 1,090 & 10 & 118 & 1,208\\
%  16th & 94 & 3,514 & 6 & 222 & 3,736\\
%  17th & 92 & 2,026 & 8 & 177 & 2,203\\
%  18th & 94 & 2,435 & 6 & 151 & 2,586\\
%  19th & 84 & 2,769 & 16 & 544 & 3,313\\
%  20th & 94 & 5,542 & 6 & 337 & 5,879\\
% \lspbottomrule
% \end{tabularx}

\fittable{
\begin{tabular}{lrrrrrrrr}
\lsptoprule
 & 13th & 14th & 15th & 16th & 17th & 18th & 19th & 20th \\
 \midrule
{\textsc{qualquier} N in \%}  & 48 & 83 & 90 & 94 & 92 & 94 & 84 & 94 \\
{\textsc{qualquier} N absolute} & 168 & 106 & 1,090 & 3,514 & 2,026 & 2,435 & 2,769 & 5,542 \\
\tablevspace
{N \textsc{qualquier} in \%}  & 52 & 17 & 10 & 6 & 8 & 6 & 16 & 6 \\
{N \textsc{qualquier} absolute}  & 185 & 21 & 118 & 222 & 177 & 151 & 544 & 337 \\
\midrule
{Total }& 353 & 127 & 1,208 & 3,736 & 2,203 & 2,586 & 3,313    & 5,879\\
\lspbottomrule
\end{tabular}
}
\caption{QN (=*ualqu* \_N*) and NQ (\_N* *ualqu*) in the whole CdE (Mar. 2020).}
\label{tab:key:22}
\end{table}

I can thus summarise that, according to the data from CdE, QN increased in time, whereas NQ decreased and there is a big jump in decrease of NQ or increase of QN in the 14th century. This big jump correlates with many other linguistic changes in this time interval of Old Spanish (see \citealt{Eberenz2009}, among others).

\ea \label{ex:key:291}
\ea \label{ex:key:291-1}  N \textsc{qualquier} (the most frequent form in the 13th c., the least frequent in the 19th c.)
\ex \label{ex:key:291-2} \textsc{qualquier} N (the less frequent form in the 13th c., the most frequent in the 19th c.)
\z
\z

In \sectref{sec:5.3}, I will explain this change in frequency of QN/NQ in terms of a linguistic change whereby \textsc{qualquier} changes its syntactic status from the specifier of a relative clause to a noun modifier (NQ) similar to an adjective (N Adj) and then to a quantificational determiner (QN) similar to \textit{algún} N.

In the next section, I will look into (\textit{otro}) NQ in the whole corpus in  more detail.

\subsection{{The} {change} {from} {(\textit{otro}}{) }{NQ} {to} {\textit{un} }{NQ }{and} {Def} {NQ}} \label{sec:5.2.2}

We will now analyse NQ more closely from the whole corpus (see \sectref{sec:5.2.1} for the frequency of NQ in total). Until the 14th century, the word order NQ appears very often in coordination with \textit{o} ‘or’ with nouns preceded by the term \textit{otro} ‘other’. The majority of the occurrences of the sequence \textit{o otro} NQ \textit{que} (120 occ. with complement clauses introduced by \textit{que} of a total of 196 occurrences of \textit{o otro} NQ, i.e. 61.22\%):

\begin{table}
\begin{tabular}{lrr}
    \lsptoprule
Century &  Frequency &    Per mil\\
    \midrule
    {ALL} & 291 & 2.91 \\
    {1200s} & 163 & 24.27 \\
    {1300s} & 8 & 3.00 \\
    {1400s} & 16 & 1.96 \\
    {1500s} & 16 & 0.94 \\
    {1600s} & 4 & 0.32 \\
    {1700s} & 20 & 2.04 \\
    {1800s} & 52 & 2.69 \\
    {1900s} & 6 & 0.26 \\
    \midrule
    {ACAD} & 2 & 0.40 \\
    {NEWS} & 1 & 0.20 \\
    {FICT} & 2 & 0.42 \\
    {ORAL} & 1 & 0.24 \\
    \midrule
     {ALL} & 2436 & 24.36 \\
    \lspbottomrule
\end{tabular}
\caption{\label{tab:distribution} Distribution of [otr* N *ualqu*]}
\end{table}

\begin{table}
\begin{tabular}{lrr}
    \lsptoprule
    {Century} & {Frequency} & {Per mil} \\
    \midrule
    {ALL} & 60 & 0.60 \\
    {1200s} & 52 & 7.74 \\
    {1300s} & 4 & 1.50 \\
    {1400s} & 2 & 0.25 \\
    {1500s} & 0 & 0.00 \\
    {1600s} & 0 & 0.00 \\
    {1700s} & 2 & 0.20 \\
    {1800s} & 0 & 0.00 \\
    {1900s} & 0 & 0.00 \\
    \midrule
    {ACAD} & 0 & 0.00 \\
    {NEWS} & 0 & 0.00 \\
    {FICT} & 0 & 0.00 \\
    {ORAL} & 0 & 0.00 \\
    \midrule
    {ALL} & 60 & 0.60 \\
    \lspbottomrule
\end{tabular}
\caption{\label{tab:distribution2} Distribution of [o otr* N *ualqu*]}
\end{table}


\ea \label{ex:key:292}
\ea \label{ex:key:292-1} 
\gll firiendo en ella con piedra, o con cuchillo o con \textit{otra} \textit{cosa} \textit{qualquier}\\
wounding in her with stone or with knife or with other thing \textsc{qualquier}\\
    \glt ‘damaging it [i.e. the Cross] with a stone, or with a knife or with some other thing whatsoever’

    (13th c., Alfonso X, \textit{Siete Partidas})
    
\ex \label{ex:key:292-2} 
\gll o \textit{otra} \textit{manera} \textit{qualquier} que le diese ayuda asabiendas que fuese semeiante o \textit{otra} \textit{cosa} \textit{qualquier} que sea\\
or another way \textsc{qualquier} that to.him gave help knowingly that it.was similar or another thing \textsc{qualquier} that is\\
    \glt ‘or [providing help to a thief in] some other way whatsoever that would be of help, knowing full well that it is similar [to the above-mentioned things] or any other thing whatever it may be’

    (13th c., Alfonso X, \textit{Siete Partidas})
    
\ex \label{ex:key:292-3} 
\gll o de \textit{otra} \textit{razon} \textit{qualquier} que sea\\
or of another reason \textsc{qualquier} that is\\
    \glt ‘or for some other reason whatever it may be’

    (13th c., Alfonso X, \textit{Siete Partidas})
\z
\z

From the 14th century onwards, there are more cases without relative clause modification, which is expected given that relative clause modification decreased in general:

\ea \label{ex:key:293}
\ea{} \label{ex:key:293-1} 
\gll […] sea debajo de la equinoccial o en \textit{otra} \textit{parte} \textit{cualquiera}» A este propósito, trae Celio Rodiginio lo de Arriano, historiador griego\\
{} is below of the equinoctial or in another part \textsc{cualquiera} to this purpose brings Celio Rodiginio what of Arriano historian Greek\\
    \glt ‘… be it under the equatorial line or in any other place” On this subject, Celio Rodiginio brings up what was said by Arriano, a Greek historian’

    (16th c., Torquemada, \textit{Jardín})
\ex \label{ex:key:293-2} 
\gll secretamente enlazado en el amor de alguna criatura, sea persona, sea \textit{otra} \textit{cosa} \textit{cualquiera}\\
secretely linked in the love of some creature is person is another thing \textsc{cualquiera}\\ 
    \glt ‘secretely bound in the love of some creature, be it a person, be it any other thing’

    (16th c., Fray Luís de Granada, \textit{Guía de pecadores})
\z
\z

All examples of \textit{otro} N Q, independently of the century, have an FCI interpretation (‘other N whichever one wants’). This observation is not surprising given that \textit{otro} N makes reference to the preceding alternatives that all belong to the same class of entities denoted by the noun N (compare English ‘fork, knife, or other utensil’), and the function of postnominal \textsc{qualquier} is to express another alternative that can have any identity or be of any type as long as it belongs to the same class of entities or types denoted by N.

Looking into the issue of whether elements other than \textit{otro} can precede the NQ sequence, we find that there are also examples of N \textsc{qualquier} with nouns that occur with a determiner. There are very few examples with the indefinite article, of the form of \textit{un/una N qual quier,-e,-a} with the separated spelling \textit{qual quier,-e,-a} in contrast to \textit{otro N qual quier} in the 14th c. (0.30 per mil vs. 37.67 per mil):

\ea \label{ex:key:294}
\ea \label{ex:key:294-1} 
\gll … Desi dixo domjngujllo que le diese el Rey \textit{vn} {omne} {qual} {quier} a que le firiese e quando le oujese ferido que se yrie para el castillo\\
{} since said Domjngujllo than to.him gave the King a man \textsc{qual} \textsc{quier} to whom to.him wounded and when to.him had.heard wounded that he would.go to the castle\\
    \glt ‘Then Dominguillo said that the King should give him a man which(ever) one wants whom he could hurt, and after hurting him, he would go to the castle.’

    (14th c., Anonymous, \textit{Veinte Reyes})
\ex \label{ex:key:294-2} 
\gll …en un logar qual quier dela posada.\\
in a place \textsc{qual} \textsc{quier} of.the inn\\
    \glt ‘…at any place in the house.’

    (13th c., Moamyn, \textit{Animalias})
\z
\z

\largerpage
The synthetic form \textsc{qualquier} in \textsc{un} N \textsc{qualquier} appears in the 16th century in its modern spelling with \textit{cu}{}- (with total of 2 occ. in this century in the whole corpus). There is no synthetic form of \textsc{un} N \textsc{qualquier} with \textit{qu-} in the whole corpus (only a separated form, as shown in 295). In these examples, \textit{cualquiera} is interpreted as an FCI due to the generic context in which \textit{cualquiera} is used (see \chapref{sec:2} on the licensing of FCIs in generic contexts):

\ea \label{ex:key:295} %ex 295
  \gll resulta difícil el destacar en \textit{una} \textit{ciencia} \textit{cualquiera}, cuánto más en la filosofía\\
results difficult to stand.out in a science \textsc{cualquiera} how.much more in the philosophy\\
  \glt ‘it turns out to be difficult to stand out in \textit{any science}, even more so in philosophy’

  (16th c., Sepúlveda, \textit{Epistolario})
\z 

\ea \label{ex:key:296} %ex 296 
  \gll que con tanto más entusiasmo nos dedicamos a \textit{un} \textit{trabajo} \textit{cualquiera}, cuanto más nos agrada el hacerlo.\\
that with so.much more enthusiasm we {dedicate ourselves} to a work \textsc{cualquiera} how.much more we like it do.it\\
  \glt ‘We dedicate ourselves to \textit{any kind of job} with the more enthusiasm, the more we like to do it.’

  (16th c., Sepúlveda, \textit{Epistolario})
\z

The use of the construction \textsc{un} N \textit{cualquiera} increases in time in the whole corpus as shown in \tabref{tab:key:syntactic}, and it has an ambiguous interpretation between the FCI interpretation (‘any’) and the evaluative interpretation (‘an ordinary’) (see \chapref{sec:3} for the ambiguity of different interpretations of \textit{cualquiera} in 20th/21st-century Spanish).

\begin{table}
% \begin{tabular}{lrrrrrrrrrrrrr}
% SECTION & ALL   & 1200s & 1300s  & 1400s & 1500s & 1600s & 1700s & 1800s & 1900s & ACAD & NEWS & FICT & ORAL \\
% FREQ    & 723    & 0    & 0      & 1     & 12    & 6     & 9     & 357    & 169  & 19   & 14    & 106    & 30    \\
% WORDS (M)& 100  & 7.9   & 3.0    & 9.7   & 19.7  & 14.8  & 11.5  & 23.1  & 22.8  & 5.0  & 5.0  & 4.8  & 4.2  \\
% PER MIL & 7.23   & 0.00 & 0.12   & 0.70  & 0.49  & 0.92  & 18.50 & 7.41  & 3.80  & 2.82 & 22.22  & 7.09 \\
% \end{tabular}

\begin{tabular}{lrrr}
\lsptoprule
Century &  Frequency &  Words (m) &  Per mil \\
\midrule
1200s &  0 &  7.9    &  0.00   \\
1300s &  0 &  3.0    &  0.00  \\
1400s &  1 &  9.7    &  0.12   \\
1500s &  12 &  19.7  &  0.70   \\
1600s &  6 &  14.8   &  0.49   \\
1700s &  9 &  11.5   &  0.92   \\
1800s &  357 &  23.1 &  18.50  \\
1900s &  169 &  22.8 &  7.41  \\
\tablevspace
ACAD &  19 &  5.0    &  3.80  \\
NEWS &  14 &  5.0    &  2.82  \\
FICT &  106 &  4.8   &  22.22 \\
ORAL &  30 &  4.2    &  7.09 \\
\midrule
ALL  &  723 & 100    &  7.23 \\
\lspbottomrule
\end{tabular}

\caption{Distribution of Indefinite Article N \textit{*ualqu*} in the whole corpus}
\label{tab:extracted7}
\end{table}

There is a single example of bare N \textit{*ualq*} under negation in the whole corpus (0 occ. from 13th to 15th c., 1 occ. in 16th c., 0 occ. in 17th, 18thc. and 19th. c.). In this example, N \textit{*ualqu*} is used under negation with the verb ‘be’ under indicative verbal mood in the present tense with an evaluative interpretation, which is expressed in the context by the contrast between N \textit{*ualqu*} and \textit{gran caudal} ‘major impact’, which triggers the interpretation of  \textit{daño cualquiera} as ‘ordinary damage’ or ‘damage of low impact’. 

\ea \label{ex:key:5-1}
\ea{}
\gll [...], el cual no es daño \textit{cualquiera}, sino de gran caudal.\\ 
{} the which \textsc{neg} is damage \textsc{cualquiera} but of great amount\\
\glt ‘which is not \textit{an ordinary damage} (EVAL), but one of major impact.’ (16th c., El Crotalón, @ 285)
\z
\z

\largerpage
The construction indefinite article N \textit{*ualq*} is also infrequent under negation per century with the exception of 19th c. where it gets more frequent (0 occ. from 13th to 16th c., 2 occ. in 17th c.(=0.16 per mil) and 18th c. (=0.10 per mil) and 10 occ.`in 19th. c. (=0.44 per mil) ). In the first example in the 17th c., \textit{*ualqu*} is used under negation with an indicative present tense verb. The examples in \REF{ex:key:297} show generic contexts with the indicative present and also have an ambiguous interpretation of ‘any’ or ‘ordinary’:

\ea \label{ex:key:297} %ex 297
\ea \label{ex:key:297-1} \gll Para pedir campo a un rey no basta \textit{un} \textit{hombre} \textit{cualquiera}, que según la ley del duelo es menester que rey sea.\\
to ask.for ground from a king \textsc{neg} is.enough a man \textsc{cualquiera} that according.to the law of.the duel is necessary that kind be\\
    \glt ‘In order to challenge a king, it is not sufficient to be \textit{an ordinary man} (EVAL), as according to the duelling law he has to be a king himself.’

‘In order to challenge a king, it is not sufficient to be \textit{just any man} (FCI), as according to the duelling law he has to be a king himself.’

    (17th c., Anonymous, \textit{El rey don Alfonso})
    
\ex \label{ex:key:297-2} \gll Quien se huelga de parecerle discreto a \textit{un} {hombre} {cualquiera}, ¿cómo piensa que trata a Dios cuando no se le da nada de no parecerle discreto?\\
who \textsc{refl} delights in seeming.to.him discreet to a man \textsc{cualquiera} how thinks that treats to God when \textsc{neg} \textsc{refl} to.him is.given nothing of \textsc{neg} seeming.to.him discreet\\
    \glt ‘Someone who likes to be seen as discreet by \textit{any man} (FCI), how does he think he treats God when he doesn’t receive anything for not being seen by him as discreet?’

    \glt ‘Someone who likes to be seen as discreet by by \textit{an ordinary man} (EVAL), how does he think he treats God when he doesn’t receive anything for not being seen by him as discreet?’

    (17th c., Zabaleta, \textit{El día de fiesta, primera parte})
    
\ex \label{ex:key:297-3} 
\gll Para instruirse en el idioma de la medicina y comer sus aforismos basta \textit{un} \textit{curso} \textit{cualquiera}.\\
to learn in the language of the medicine and eat its aphorisms is.enough a course \textsc{cualquiera}\\
    \glt ‘in order to learn the language of medicine and to eat its aphorisms, any course (FCI)/an ordinary course (EVAL) is sufficient.’ (1732, \textit{Vida}, Torres Villarroel)

    (17th c., Zabaleta, \textit{El día de fiesta, primera parte})
\z
\z 

From this time onwards, \textsc{un} N \textit{cualquiera} is used as EVAL or FCI depending on the context as in present-day Spanish (see \chapref{sec:3}). In the following example, \textit{un alojamento cualquiera} can be interpreted as FCI ‘any accomodation’ or as ‘a normal accomodation’ excluding ‘abnormal’ or ‘unlivable’ accomodations:

\ea \label{ex:key:298} %ex 298
    \gll En esta perplejidad escribo al administrador del conde de Peñalba, encargándole que me busque \textit{un} \textit{alojamiento} \textit{cualquiera} de acuerdo con usted, ....\\
in this perplexity write to.the administrator of=the count of Peñalba entrusting=him that to.me find a accomodation \textsc{cualquiera} in accordance with you\\
    \glt ‘I’m writing in this perplexity to the administrator of the Count of Peñalba, charging him to find me some accomodation in accordance with you, ...’

    (18th c., Jovellanos, \textit{Correspondencia})
\z

The first examples of \textsc{un} N \textit{cualquiera} in affirmative predicative contexts appeared in the 19th century:

\ea \label{ex:key:299} %ex 299
    \gll Es \textit{un} \textit{contrato} \textit{cualquiera}. vos necesitáis dinero, yo os necesito, ....\\
is a contract \textsc{cualquiera} you need money I you need\\
    \glt ‘It’s a usual contract: you need money, I need you, …’

    (19th c., Zorrilla, \textit{Entre clérigos y diablos})
\z

The use of \textit{N cualquiera} with definite articles or demonstratives such \textit{la, el, esta, este, aquel} N (Def N \textit{cualquiera}) is less frequent than the use of \textit{un/otro} N \textit{cualquiera} (less than 5\% of all occurrences) and also appears only in the 19th century (see \chapref{sec:3} for present-day Spanish). An explanation of the late appearance of Def N \textit{cualquiera} will be given in \sectref{sec:5.3}. All examples of Def \textit{N cualquiera} have an evaluative interpretation of ‘common/ordinary’ (see also \chapref{sec:3} for Modern Spanish):

\ea \label{ex:key:300} %ex 300
  \gll La cuestión empezó \textit{ese} \textit{día} \textit{cualquiera} del siglo XVI, en que don Felipe II decretó que los españoles no podrían estudiar en Universidades extranjeras.\\
the question started that day \textsc{cualquiera} of.the century XVI in which Don Felipe II decreed that the Spaniards \textsc{neg} could study in universities foreign\\
  \glt ‘the problem began on this ordinary day of the 16th century, when Don Felipe decreed, that Spaniards could not study at foreign universities.’

  (20th c., Ballester, \textit{Torre Del Aire})
\z

Given these observations, it is now possible to explain the rise of NQ in the 18th and 19th centuries observed in \sectref{sec:5.2.1.2}. The introduction of the new semantic meaning of \textsc{un} NQ ‘an ordinary N’ correlates with the introduction of new syntactic constructions: the use of \textsc{un} NQ under affirmative predicative contexts and the use of NQ under definite articles. These new forms and the new meaning have triggered the increase in frequency of NQ in the 18th and 19th centuries. 
As the construction (\textit{otro}) NQ practically disappeared in the 19th century and Q(\textit{otro})N was used instead (\tabref{tab:key:19}), this explains the trend of the general fall of NQ after the 14th century (see \sectref{sec:5.2.1.2}).

As noted above, there is a renewed decrease of postnominal \textsc{qualquier} from the 18th to 19th century, although the relative frequency of postnominal \textsc{cualquiera} in 19th c. is still higher than the relative frequency of postnominal \textsc{qualquier} before the 18th centuery (see \tabref{tab:key:12} in my sample and \tabref{tab:key:syntactic} in the whole corpus). If the renewed decrease were an effect of random sampling in my sample in  \tabref{tab:key:12}, we would not see the same effect appearing in the whole corpus (see \tabref{tab:key:syntactic}). This effect might be triggered by the differences in the selection of text genre in 18th and 19th c. The 19th c. has a more nuanced text genre constellation: academic texts, news, fiction and texts representing oral speech. Indeed, \tabref{tab:extracted3} shows that the highest distribution of postnominal \textsc{qualquier} is in the genres fiction and oral texts. I thus expect the decrease to disappear or to be less strong if the genre is more controlled in the diachronic study. How genre affects diachronic change of \textsc{qualquier} will be tested in future research.

To sum up the development of N \textit{cualquiera} (bold letters represent higher frequency). It represents a change in different constructions with N \textit{cualquiera.} It shows that the construction \textit{otro/otra} \textsc{qualquier} (\textit{que} Verb\textsubscript{Subj.}) was first used in the 13th century until the 15th. The construction \textsc{un} N \textit{cualquiera} appeared in the 15th century and the older construction disappeared. From the 19th century onwards, \textsc{un} N \textit{cualquiera} was used in predicative position. Fewer data of N \textit{cualquiera} contain a definite or demonstrative \textit{la}/\textit{el}/\textit{este}/\textit{esta} N \textit{cualquiera}:

% \ea \label{ex:key:301} %ex 301
%   \& \textit{otro,a} \textsc{qualquier} \textit{que} Verb\textsubscript{Subj.} (13th–15th c.)

%     (Verb\textsubscript{Ind.}) \textsc{un} N \textit{cualquiera} (Verb\textsubscript{Ind.}) (from 15th c. until today) $\rightarrow$

%    {pedicative} {verb\textsubscript{Ind.}} {\textsc{un}} {N} {\textit{cualquiera}} /\textit{la,el,este,esta} N \textit{cualquiera} (from 19th c. until today)
% \z

\begin{table}
\begin{tabularx}{\textwidth}{@{}lX@{}}
\lsptoprule
{13th–15th c.} & \textit{otro,a} \textsc{qualquier} \textit{que} Verb\textsubscript{Subj.} \\
{15th c.–today} & (Verb\textsubscript{Ind.}) \textsc{un} N \textit{cualquiera} (Verb\textsubscript{Ind.}) \\
{19th c.–today} & \textit{la, el, este, esta} N \textit{cualquiera} \\
                       & [predicative] Verb\textsubscript{Ind.} \textsc{un} N \textit{cualquiera} \\
\lspbottomrule
\end{tabularx}
\caption{Historical development of \textit{cualquiera} \label{ex:key:301}}
\end{table}


The distribution in \tabref{ex:key:301} co-occurs with a semantic distribution given the observation that \textit{otro} NQ can only have FCI due to preceding alternatives, whereas \textit{un} NQ is ambiguous between two interpretations (FCI/EVAL) depending on the verbal mood (+/– modal) and Def NQ has only the interpretation of EVAL:


\ea \label{ex:key:302} %ex 302
  \textit{otro,a} N \textsc{qualquier} \textsc{que} ….FCI (13th–15th c.)

      \textsc{un} N \textit{cualquiera} FCI/\textbf{EVAL} (from the 16th c. until today)

   \textit{la,el,este,esta} N \textit{cualquiera} \textbf{EVAL} (from the 19th c. until today)
\z

As I will argue in \sectref{sec:5.3}, the preference for \textsc{un} N \textit{cualquiera} over \textit{(otro)} N \textit{qual quiera} rising and later the appearance of Def N \textit{cualquiera} is the result of the reanalysis of \textit{cualquiera} from an item introducing a relative clause with FCI into a synthetic nominal modifier with different lexical meanings (FCI and EVAL).
 The relevant diachronic steps of postnominal \textit{cualquiera} in European Spanish that have triggered EVAL are summarised as follows:

  \begin{enumerate}[label=\arabic*., leftmargin=2em]
    \item Use of postnominal \textit{cualquiera} often after \textit{otro ‘other’ + N} with a complement clause after \textit{quier} expressed by \textit{que ‘that’} and a selected subjunctive (e.g. \textit{qual quier que sea}) with modal interpretation ‘whichever one wants that it may be’ in 13th century Old Spanish:
    \ea 
        \gll \textit{e.g.} \textit{o} \textit{de} \textit{otra} \textit{razon} \textit{qualquier} \textit{que} \textit{sea}\\
        {} or of another reason \textsc{qualquier} that may.be\\
        \glt ‘or for some other reason whichever one wants that it may be’ 
        
        (13th c., Alfonso X, \textit{Siete Partidas})
    \z 
    
    Use of nouns between \textit{qual} and \textit{quier}, which show that \textit{qual quier} could be analysed compositionally in 13th century \citep{Pato2012}:
    \ea
        \gll \textit{e.g.} {[...]} {en} {qual} {manera} {quier} {que} {sean} \\
        {} {} \textit{in} \textit{which} \textit{manner} \textit{wants} \textit{that} \textit{be-subj} \\
        \glt ‘in whatever manner’ 
        
        (13th c., Alfonso X, \textit{Libro de las cruces})
    \z

    \item Recategorisation of postnominal \textit{cualquiera} as one word, the compositional meaning of the relative clause (‘which one wants’) is no longer expressed as a sentence but as a single lexical meaning ‘any’). Empirical evidence: 
    \begin{enumerate}[label=\alph*), leftmargin=1.5em]
        \item disappearance of \textit{qual NP quier}, 
        \item orthographic non-separated form \textit{qualquier} increases, 
        \item decrease of complement clause after \textit{quier} and selected subjunctive: \textit{qual quier que ….}, 
        \item increasing use of postnominal \textit{cualquiera} with an indefinite N as in \textit{\textsc{un} N cualquiera}.
    \end{enumerate}

    \item Use of \textit{cualquiera} in predicative and postnominal position under negation (e.g. \textit{no es un día cualquiera, es un día especial} ‘it’s not an ordinary day, it’s a special day’) (in 16th c. in European Spanish).
    
    \item Increase of copular predicates without negation in present and past tense (e.g. \textit{es un día cualquiera} ‘it’s an ordinary day’) (from 18th c. in European Spanish).
    
    \item Extension of the use of \textit{\textsc{un} N+ cualquiera} to definite articles + N \textit{cualquiera} and quantified nouns as in \textit{la/mucha gente cualquiera} ‘the/many ordinary people’ (in European Spanish in the 20th century).

\end{enumerate}

\subsection{Summing up and evaluating the results}\label{sec:5.2.3}

It has been shown that the frequency of some variables increased , such as postverbal \textit{cualquiera} (VQ), sentence types other than relative clauses, the verb mood indicative and QN. Moreover, \textsc{un} N \textsc{cualquiera} came into existence in later centuries (\tabref{tab:key:syntactic}).


\begin{table}
\caption{Syntactic and semantic changes related to \textsc{qualquier}}
\begin{tabularx}{\textwidth}{X}
\lsptoprule
{\bfseries {In early centuries (12th–13th c.)} }\\
\midrule
  NQ is more frequent than QN.

  Q \textit{que} is more frequent than other sentence types.

  Subjunctive is more frequent than indicative.

  \textsc{un} N \textsc{qualquier} with indicative did not exist yet.

  \textsc{qualquier} is only interpretable as FCI.\\

\tablevspace
  {\bfseries {In later centuries (19th–20th c.)}}\\
\midrule
  QN is more frequent than NQ.

  Q \textit{que} is less frequent than other sentence types.

  Indicative is more frequent than subjunctive.

  \textsc{un} N \textsc{qualquier} with indicative verbal mood appears.

  \textsc{qualquier} has different meanings (FCI, EVAL).\\
\lspbottomrule
\end{tabularx}
\label{tab:key:syntactic}
\end{table}

The main results suggest that \textsc{qualquier} has undergone some changes which correlate with syntatic position and verbal mood/type and (in)definiteness of the NP.

As was shown in \sectref{sec:2.6}, no quantitative analysis had been done so far on the variables described in \sectref{sec:5.2}, namely the syntactic function of \textsc{qualquier}, its relation to the verb (position and mood), and its relation to the sentence mood. I have also shown when exactly new constructions \textsc{un}/Def N \textsc{qualquier} appear for the first time, and I have quantified them. The construction \textsc{un} N \textsc{qualquier} appears already in the 13th century, but it is used very infrequently in contrast to \textit{otro} N \textsc{qualquier}. In later centuries, \textsc{un} N \textsc{qualquier} starts being used more frequently (from the 18th c. onwards). The construction Def N \textsc{qualquier} is a recent phenomenon (from the 19th c. onwards) which is still less frequent than \textsc{un} N \textsc{qualquier.} I have identified the time interval when the meanings of \textsc{qualquier} started to change, namely from the 16th century.

I will use these results to investigate the meaning changes of \textsc{qualquier}. Before turing to the semantic side of linguistic changes, I will check whether the described linguistic changes as represented in my sample exist in the whole corpus if we ignore factors like the philological reliability of certain texts and manual tagging of examples.


\subsection{Quantitative results in the whole corpus}\label{sec:5.2.4-0} %new label added, section is not in manuscript

The changes related to \textsc{qualquier} as represented in the whole corpus can be summarised as follows.

\begin{itemize}

\item The use of the orthographic variant \textsc{qualquier} decreases and the use of \textsc{cualquier} increases (see \tabref{tab:extracted1} and \tabref{tab:extracted2})
\item Prenominal \textsc{qualquier} (QN) increases (see \tabref{tab:extracted3}) and NQ without an article [bare N Q] decreases (see \tabref{tab:extracted4}).
\item The use of bare \textsc{qualquier} (bare Q) without nominal modification and the use of \textsc{qualquier} \textit{que} (Q \textit{que}) decreases (see \tabref{tab:extracted5}).
\item the postverbal position increases (VQ) increases (see \tabref{tab:extracted6}, \tabref{tab:extracted7-2} and \tabref{tab:extracted8})
\item the use of \textsc{qualquier} under predicative verbs like \textit{be} increases (see \tabref{tab:extracted})
\item the use of \textsc{qualquier} indicative verbal mood increases (see \tabref{tab:extracted9})
\item the use of postnominal \textsc{qualquier} in the 20th c. is clearly related to the genre (fictional and oral texts) (see the higher distribution of fictional and oral texts in comparison to other genres in \tabref{tab:extracted6}, \tabref{tab:extracted6}, \tabref{tab:extracted7-2}, \tabref{tab:extracted8})

\end{itemize}

%}
%\documentclass{article}
%\usepackage{graphicx} % For including graphics
%\usepackage{array}    % For extending table functionalities


\begin{table}
\caption{Distribution of \textit{qualqu*} in the whole corpus}
\label{tab:extracted1}
% \begin{tabular}{lrrrrrrrrrrrrr}
% SECTION & ALL  & 1200s & 1300s & 1400s & 1500s & 1600s & 1700s & 1800s & 1900s & ACAD & NEWS & FICT & ORAL \\
% FREQ    & 4392 & 1008  & 213   & 1688  & 782   & 234   & 436   & 27    & 2     & 0    & 0    & 0    & 2    \\
% WORDS (M)& 100  & 7.9   & 3.0   & 9.7   & 19.7  & 14.8  & 11.5  & 23.1  & 22.8  & 5.0  & 5.0  & 4.8  & 4.2  \\
% PER MIL & 43.92& 150.10& 79.79 & 206.83& 45.91 & 18.95 & 44.41 & 1.40  & 0.09  & 0.00 & 0.00 & 0.00 & 0.47 \\
% \end{tabular}

\begin{tabular}{lrrrr}
\lsptoprule
Century &  Frequency &  Words (m) &  Per mil\\
\midrule
1200s &  1008 &  7.9 &  150.10 \\
1300s &  213 &  3.0 &  79.79 \\
1400s &  1688 &  9.7 &  206.83 \\
1500s &  782 &  19.7 &  45.91 \\
1600s &  234 &  14.8 &  18.95 \\
1700s &  436 &  11.5 &  44.41\\
1800s &  27 &  23.1 &  1.40 \\
1900s &  2 &  22.8 &  0.09 \\
\tablevspace
ACAD &  0 &  5.0 &  0.00 \\
NEWS &  0 &  5.0 &  0.00 \\
FICT &  0 &  4.8 &  0.00 \\
ORAL  &  2     &  4.2   &  0.47 \\
\midrule
 {ALL} &  4392 &  100 &  43.92 \\
\lspbottomrule
\end{tabular}
\end{table}

\begin{table}
\caption{Distribution of \textit{cualqu*} in the whole corpus}
\label{tab:extracted2}
% \begin{tabular}{lrrrrrrrrrrrrr}
% \textbf{SECTION} & \textbf{ALL}  & \textbf{1200s} & \textbf{1300s} & \textbf{1400s} & \textbf{1500s} & \textbf{1600s} & \textbf{1700s} & \textbf{1800s} & \textbf{1900s} & \textbf{ACAD} & \textbf{NEWS} & \textbf{FICT} & \textbf{ORAL} \\
% \textbf{FREQ}    & 32364 & 2  & 3   & 171  & 4799 & 2691 & 3611 & 5611 & 7738 & 1611 & 1787 & 2176 & 2164 \\
% \textbf{WORDS (M)} & 100   & 7.9  & 3.0   & 9.7   & 19.7  & 14.8  & 11.5  & 23.1  & 22.8  & 5.0  & 5.0  & 4.8  & 4.2  \\
% \textbf{PER MIL} & 323.64 & 0.30 & 1.12 & 20.95 & 281.73 & 217.93 & 367.84 & 290.77 & 339.05 & 322.20 & 360.01 & 456.20 & 511.21 \\
% \end{tabular}
\begin{tabular}{lrrrr}
\lsptoprule
Century &  Frequency &  Words (m) &  Per mil\\
\midrule
{1200s} & 2 & 7.9 & 0.30 \\
{1300s} & 3 & 3.0 & 1.12 \\
{1400s} & 171 & 9.7 & 20.95 \\
{1500s} & 4799 & 19.7 & 281.73 \\
{1600s} & 2691 & 14.8 & 217.93 \\
{1700s} & 3611 & 11.5 & 367.84 \\
{1800s} & 5611 & 23.1 & 290.77 \\
{1900s} & 7738 & 22.8 & 339.05 \\
\tablevspace
{ACAD} & 1611 & 5.0 & 322.20 \\
{NEWS} & 1787 & 5.0 & 360.01 \\
{FICT} & 2176 & 4.8 & 456.20 \\
{ORAL} & 2164 &  4.2 & 511.21 \\
\midrule
{ALL} & 32364 & 100 & 323.64 \\
\lspbottomrule
\end{tabular}
\end{table}

\begin{table}
\caption{Distribution of N \textsc{*ualqu*} in the whole corpus}
\label{tab:extracted3}
% \begin{tabular}{lrrrrrrrrrrrrr}
% \textbf{SECTION} & \textbf{ALL}  & \textbf{1200s} & \textbf{1300s} & \textbf{1400s} & \textbf{1500s} & \textbf{1600s} & \textbf{1700s} & \textbf{1800s} & \textbf{1900s} & \textbf{ACAD} & \textbf{NEWS} & \textbf{FICT} & \textbf{ORAL} \\
% \textbf{FREQ}    & 2436 & 501  & 40   & 176  & 330 & 175 & 335 & 527 & 176 & 30 & 48 & 43 & 55 \\
% \textbf{WORDS (M)} & 100   & 7.9  & 3.0   & 9.7   & 19.7  & 14.8  & 11.5  & 23.1  & 22.8  & 5.0  & 5.0  & 4.8  & 4.2  \\
% \textbf{PER MIL} & 24.36 & 74.60 & 14.98 & 21.56 & 19.37 & 14.17 & 34.13 & 27.31 & 7.71 & 6.00 & 9.67 & 9.01 & 12.99 \\
% \end{tabular}
\begin{tabular}{lrrr}
\lsptoprule
Century &  Frequency &  Words (m) &  Per mil\\
\midrule
{1200s} &  501 &  7.9 &  74.60   \\
{1300s} &  40 &  3.0 &  14.98   \\
{1400s} &  176 &  9.7 &  21.56   \\
{1500s} &  330 &  19.7 &  19.37   \\
{1600s} &  175 &  14.8 &  14.17   \\
{1700s} &  335 &  11.5 &  34.13   \\
{1800s} &  527 &  23.1 &  27.31   \\
{1900s} &  176 &  22.8 &  7.71   \\
\tablevspace
{ACAD} &  30 &  5.0 &  6.00   \\
{NEWS} &  48 &  5.0 &  9.67   \\
{FICT} &  43 &  4.8 &  9.01   \\
{ORAL}  &  55 &  4.2   &  12.99 \\
\midrule
{ALL} &  2436 &  100 &  24.36    \\
\lspbottomrule
\end{tabular}
\end{table}

\begin{table}
\caption{Distribution of \textit{*ualqu*} N in the whole corpus}
\label{tab:extracted4}
% \begin{tabular}{lrrrrrrrrrrrrr}
% SECTION & ALL  & 1200s & 1300s & 1400s & 1500s & 1600s & 1700s & 1800s & 1900s & ACAD & NEWS & FICT & ORAL \\
% FREQ    & 23192 & 168  & 106   & 1090  & 3514   & 2026   & 2435   & 2769    & 5542     & 1203    & 1386    & 1412    & 1541    \\
% WORDS (M)& 100  & 7.9   & 3.0   & 9.7   & 19.7  & 14.8  & 11.5  & 23.1  & 22.8  & 5.0  & 5.0  & 4.8  & 4.2  \\
% PER MIL & 231.92 & 25.02& 39.71 & 133.56& 206.29 & 164.07 & 248.05 & 143.49  & 242.83  & 240.60 & 279.22 & 296.02 & 364.04 \\
% \end{tabular}
\begin{tabular}{lrrr}
\lsptoprule
Century &  Frequency &  Words (m) &  Per mil\\
\midrule
 1200s &  168 & 7.9 & 25.02 \\
 1300s &  106 & 3.0 & 39.71 \\
 1400s &  1090 & 9.7 &  133.56 \\
 1500s & 3514 &  19.7 & 206.29 \\
 1600s & 2026 &  14.8 &  164.07 \\
 1700s & 2435 &  11.5 & 248.05 \\
 1800s & 2769 & 23.1 &  143.49 \\
 1900s & 5542 & 22.8 & 242.83 \\
 \tablevspace
 ACAD &  1203 & 5.0 & 240.60 \\
 NEWS &  1386 & 5.0 & 279.22 \\
 FICT &  1412 & 4.8 & 296.02 \\
 ORAL &  1541  & 4.2  & 364.04 \\
 \midrule
ALL & 23192 &  100 & 231.92 \\
\lspbottomrule
\end{tabular}
\end{table}

\begin{table}
% \begin{tabular}{lrrrrrrrrrrrrr}
% SECTION & ALL  & 1200s & 1300s & 1400s & 1500s & 1600s & 1700s & 1800s & 1900s & ACAD & NEWS & FICT & ORAL \\
% FREQ    & 2436 & 501  & 40   & 176  & 330   & 175   & 335   & 527    & 176     & 30    & 48    & 43    & 55    \\
% WORDS (M)& 100  & 7.9   & 3.0   & 9.7   & 19.7  & 14.8  & 11.5  & 23.1  & 22.8  & 5.0  & 5.0  & 4.8  & 4.2  \\
% PER MIL & 24.36 & 74.60& 14.98 & 21.56& 19.37 & 14.17 & 34.13 & 27.31  & 7.71  & 6.00 & 9.67 & 9.01 & 12.99 \\
% \end{tabular}
\begin{tabular}{lrrr}
\lsptoprule
Century &  Frequency &  Words (m) &  Per mil\\
\midrule
1200s& 501& 7.9& 74.60\\
1300s& 40& 3.0& 14.98\\
1400s& 176& 9.7& 21.56\\
1500s& 330& 19.7& 19.37\\
1600s& 175& 14.8& 14.17\\
1700s& 335& 11.5& 34.13\\
1800s& 527& 23.1& 27.31\\
1900s& 176& 22.8& 7.71\\
\tablevspace
ACAD& 30& 5.0& 6.00\\
NEWS& 48& 5.0& 9.67\\
FICT& 43& 4.8& 9.01\\
ORAL & 55& 4.2  & 12.99 \\
\midrule
ALL& 2436& 100& 24.36 \\
\lspbottomrule
\end{tabular}
\caption{Distribution of \textit{*ualqu*} que in the whole corpus}
\label{tab:extracted5}
\end{table}

\begin{table}
% \begin{tabular}{lrrrrrrrrrrrrr}
% SECTION  & ALL  & 1200s & 1300s & 1400s & 1500s & 1600s & 1700s & 1800s  & 1900s & ACAD & NEWS & FICT  & ORAL \\
% FREQ     & 900  & 2     & 1     & 18    & 40    & 36    & 29    & 380    & 197   & 21   & 19   & 119   & 38    \\
% WORDS (M)& 100  & 7.9   & 3.0   & 9.7   & 19.7  & 14.8  & 11.5  & 23.1   & 22.8  & 5.0  & 5.0  & 4.8   & 4.2  \\
% PER MIL  & 9.00 & 0.30  & 0.37  & 2.21  & 2.35  & 2.92  & 2.95  & 19.69  & 8.63  & 4.20 & 3.83 & 24.95 & 8.98 \\
% \end{tabular}

\begin{tabular}{lrrr}
\lsptoprule
Century &  Frequency &  Words (m) &  Per mil\\
\midrule
  1200s& 2& 7.9& 0.30\\
  1300s& 1& 3.0& 0.37\\
  1400s& 18& 9.7& 2.21\\
  1500s& 40& 19.7& 2.35\\
  1600s& 36& 14.8& 2.92\\
  1700s& 29& 11.5& 2.95\\
  1800s& 380& 23.1& 19.69\\
  1900s& 197& 22.8& 8.63\\
  \tablevspace
  ACAD& 21& 5.0& 4.20\\
  NEWS& 19& 5.0& 3.83\\
  FICT& 119& 4.8& 24.95\\
  ORAL & 38    & 4.2  & 8.98 \\
  \midrule
  ALL& 900& 100& 9.00\\
\lspbottomrule
\end{tabular}

\caption{Distribution of Article N \textit{*ualqu*} in the whole corpus}
\label{tab:extracted6}
\end{table}

\begin{table}
% \begin{tabular}{lrrrrrrrrrrrrr}
% SECTION  & ALL    & 1200s & 1300s & 1400s & 1500s & 1600s & 1700s & 1800s & 1900s & ACAD & NEWS   & FICT & ORAL \\
% FREQ     & 723    & 0     & 0     & 1     & 12    & 6     & 9     & 357   & 169   & 19   & 14     & 106    & 30    \\
% WORDS (M)& 100    & 7.9   & 3.0   & 9.7   & 19.7  & 14.8  & 11.5  & 23.1  & 22.8  & 5.0  & 5.0    & 4.8  & 4.2  \\
% PER MIL  & 7.23   & 0.00  & 0.12  & 0.70  & 0.49  & 0.92  & 18.50 & 7.41  & 3.80  & 2.82 & 22.22  & 7.09 & ??? \\
% \end{tabular}
\begin{tabular}{lrrr}
\lsptoprule
Century &  Frequency &  Words (m) &  Per mil\\
  \midrule
  1200s& 0& 7.9& 0.00\\
  1300s& 0& 3.0& 0.12\\
  1400s& 1& 9.7& 0.70\\
  1500s& 12& 19.7& 0.49\\
  1600s& 6& 14.8& 0.92\\
  1700s& 9& 11.5& 18.50\\
  1800s& 357& 23.1& 7.41\\
  1900s& 169& 22.8& 3.80\\
  \tablevspace
  ACAD& 19& 5.0& 2.82\\
  NEWS& 14& 5.0& 22.22\\
  FICT& 106& 4.8& 7.09\\
  ORAL & 30    & 4.2  & 7.09 \\
  \midrule
  ALL& 723& 100& 7.23\\
  \lspbottomrule
\end{tabular}

\caption{Distribution of Indefinite Article N \textit{*ualqu*} in the whole corpus}
\label{tab:extracted7-2}
\end{table}

\begin{table}
% \begin{tabular}{lrrrrrrrrrrrrr}
% SECTION & ALL  & 1200s & 1300s & 1400s & 1500s & 1600s & 1700s & 1800s & 1900s & ACAD & NEWS & FICT & ORAL \\
% FREQ    & 117 & 0  & 0   & 0  & 1 & 1  & 2   & 43   & 35    & 3     & 2    & 22    & 8    \\
% WORDS (M)& 100  & 7.9   & 3.0   & 9.7   & 19.7  & 14.8  & 11.5  & 23.1  & 22.8  & 5.0  & 5.0  & 4.8  & 4.2  \\
% PER MIL & 1.17   & 0.00   & 0.00   & 0.00  & 0.06  & 0.08  & 0.20  & 2.23 & 1.53  & 0.60  & 0.40  & 4.61  & 1.89 \\
% \end{tabular}
\begin{tabular}{lrrr}
\lsptoprule
Century &  Frequency &  Words (m) &  Per mil\\
  \midrule
  1200s& 0& 7.9& 0.00\\
  1300s& 0& 3.0& 0.00\\
  1400s& 0& 9.7& 0.00\\
  1500s& 1& 19.7& 0.06\\
  1600s& 1& 14.8& 0.08\\
  1700s& 2& 11.5& 0.20\\
  1800s& 43& 23.1& 2.23\\
  1900s& 35& 22.8& 1.53\\
  \tablevspace
  ACAD& 3& 5.0& 0.60\\
  NEWS& 2& 5.0& 0.40\\
  FICT& 22& 4.8& 4.61\\
  ORAL & 8    & 4.2  & 1.89 \\
  \midrule
  ALL& 117& 100& 1.17\\
\lspbottomrule
\end{tabular}

\caption{Distribution of Verb Indefinite Article N \textit{*ualqu*} in the whole corpus}
\label{tab:extracted8}
\end{table}

\begin{table}
% \begin{tabular}{lrrrrrrrrrrrrr}
% SECTION & ALL  & 1200s & 1300s & 1400s & 1500s & 1600s & 1700s & 1800s & 1900s & ACAD & NEWS & FICT & ORAL \\
% FREQ    & 41 & 0  & 0   & 0  & 0 & 1  & 2   & 14   & 12    & 1     & 2    & 3    & 6    \\
% WORDS (M)& 100  & 7.9   & 3.0   & 9.7   & 19.7  & 14.8  & 11.5  & 23.1  & 22.8  & 5.0  & 5.0  & 4.8  & 4.2  \\
% PER MIL &  0.41   & 0.00   & 0.00   & 0.00  & 0.00  & 0.08  & 0.20  & 0.73 & 0.53  & 0.20  & 0.40  & 0.63  & 1.42 \\
% \end{tabular}

\begin{tabular}{lrrr}
\lsptoprule
Century &  Frequency &  Words (m) &  Per mil\\
\midrule
1200s& 0& 7.9& 0.00\\
1300s& 0& 3.0& 0.00\\
1400s& 0& 9.7& 0.00\\
1500s& 0& 19.7& 0.00\\
1600s& 1& 14.8& 0.08\\
1700s& 2& 11.5& 0.20\\
1800s& 14& 23.1& 0.73\\
1900s& 12& 22.8& 0.53\\
ACAD& 1& 5.0& 0.20\\
NEWS& 2& 5.0& 0.40\\
FICT& 3& 4.8& 0.63\\
ORAL & 6    & 4.2  & 1.42 \\
\midrule
ALL& 41& 100& 0.41\\
\lspbottomrule
\end{tabular}
\caption{Distribution of Present Tense Verb Indefinite Article N \textit{*ualqu*} in the whole corpus}
\label{tab:extracted9}
\end{table}

\begin{table}
% \begin{tabular}{lrrrrrrrrrrrrrr}
% SECTION & ALL  & 1200s & 1300s & 1400s & 1500s & 1600s & 1700s & 1800s & 1900s & ACAD & NEWS & FICT & ORAL \\
% FREQ    & 6 & 0  & 0   & 0  & 0 & 0  & 0   & 0   & 3    & 0     & 0    & 3    & 0    \\
% WORDS (M)& 100  & 7.9   & 3.0   & 9.7   & 19.7  & 14.8  & 11.5  & 23.1  & 22.8  & 5.0  & 5.0  & 4.8  & 4.2  \\
% PER MIL & & 0.06   & 0.00   & 0.00   & 0.00  & 0.00  & 0.00  & 0.00  & 0.00 & 0.13  & 0.00  & 0.00  & 0.63  & 0.00 \\
% \end{tabular}

\begin{tabular}{lrrr}
\lsptoprule
Century &  Frequency &  Words (m) &  Per mil\\
  \midrule
  1200s& 0& 7.9& 0.00\\
  1300s& 0& 3.0& 0.00\\
  1400s& 0& 9.7& 0.00\\
  1500s& 0& 19.7& 0.00\\
  1600s& 0& 14.8& 0.00\\
  1700s& 0& 11.5& 0.00\\
  1800s& 0& 23.1& 0.00\\
  1900s& 3& 22.8& 0.13\\
  \tablevspace
  ACAD& 0& 5.0& 0.00\\
  NEWS& 0& 5.0& 0.00\\
  FICT& 3& 4.8& 0.63\\
  ORAL & 0    & 4.2  & 0.00\\
  \midrule
  ALL& 6& 100& 0.06 \\
  \lspbottomrule
\end{tabular}

\caption{Distribution of \textit{Indefinido} Indefinite Article N \textit{*ualqu*} in the whole corpus}
\label{tab:extracted10}
\end{table}

\begin{table}
% \begin{tabular}{lrrrrrrrrrrrrrr}
% SECTION & ALL  & 1200s & 1300s & 1400s & 1500s & 1600s & 1700s & 1800s & 1900s & ACAD & NEWS & FICT & ORAL \\
% FREQ    & 52 & 0  & 0   & 0  & 1 & 0  & 0   & 5   & 23    & 2     & 4    & 10    & 7    \\
% WORDS (M)& 100  & 7.9   & 3.0   & 9.7   & 19.7  & 14.8  & 11.5  & 23.1  & 22.8  & 5.0  & 5.0  & 4.8  & 4.2  \\
% PER MIL & & 0.52   & 0.00   & 0.00   & 0.00  & 0.06  & 0.00  & 0.00  & 0.26 & 1.01  & 0.40  & 0.81  & 2.10  & 1.65 \\
% \end{tabular}

\begin{tabular}{lrrr}
\lsptoprule
Century &  Frequency &  Words (m) &  Per mil\\
  \midrule
  1200s& 0& 7.9& 0.00\\
  1300s& 0& 3.0& 0.00\\
  1400s& 0& 9.7& 0.00\\
  1500s& 1& 19.7& 0.06\\
  1600s& 0& 14.8& 0.00\\
  1700s& 0& 11.5& 0.00\\
  1800s& 5& 23.1& 0.26\\
  1900s& 23& 22.8& 1.0\\
  \tablevspace
  ACAD& 2& 5.0& 0.40\\
  NEWS& 4& 5.0& 0.81\\
  FICT& 10& 4.8& 2.10\\
  ORAL & 7    & 4.2  & 1.65\\
  \midrule
  ALL& 52& 100& 0.52\\
  \lspbottomrule
\end{tabular}

\caption{Distribution of the verb \textit{be} Indefinite Article N \textit{*ualqu*} in the whole corpus}
\label{tab:extracted}
\end{table}

To sum up, the distribution of the linguistic variables considered in the whole corpus matches the distribution of the same variables in my sample. 

\section{The semantic side of the grammaticalisation}\label{sec:5.3}

In this section, I will spell out the change from a relative clause into a quantificational pronoun and determiner from the semantic point of view. The change into a postnominal modifier will be discussed as well.

\subsection{{The} {prenominal} {determiner}} \label{sec:5.3.1}

I assume that the determiner/pronoun \textit{qualquier} has its origin in the relative clause and that it is the quantificational interpretation of a free relative clause such as \textit{qual quier(e) que} that triggers the reanalysis into a quantificational determiner. I will first illustrate this assumption with the interpretation of a relative clause with a free choice interpretation in Modern Spanish and then apply the semantic interpretation to Old Spanish. As \REF{ex:key:305} shows, free relative clauses can be replaced by a universal quantificational pronoun ‘everyone’:

\ea \label{ex:key:305} %ex 305
 \ea 
 \gll Puedes invitar a quien tú quieras.\\
can invite to whom you want {} \\
    \glt ‘You can invite whom you want (to invite).’

\ex \gll Puedes invitar a todos (los que tú quieras invitar)\\
can invite to all (the that you want invite)\\
\glt ‘You can invite everyone (you want to invite).’  
\z
\z 

I adopt the denotation of English Free Relatives (FRs) in English (see \citealt{Simik2018}) to (FRs) in Spanish. The truth-conditions of \REF{ex:key:305} are provided in \REF{ex:key:306}.

\ea \label{ex:key:306} %ex 306
\ea{} \label{ex:key:306-1}  ⟦quien tu quieras invitar⟧ = λP.${\forall}$x [want′(x)(you′) → P(x)]

(the set of properties that all people that you want to invite)
\ex{} \label{ex:key:306-2} ⟦(305)⟧ = 1 iff

{${\forall}$x [want′ (x)(you′)(invite′)] → can (invite′(x)(you))}

{For all people you want to invite, you can invite them.}
\z
\z

Free relatives can thus be analysed as DPs which can be modelled by the insertion of a covert DP with the definite interpretation ‘the NP that’ or universal interpretation (‘everyone’) depending on the verb mood above the relative clause structure; free relatives embedded under modal verbs usually have a universal interpretation:

\ea \label{ex:key:307} %ex 307
{  Puedes invitar [\textsubscript{DP} a todos [\textsubscript{CP} a}\footnote{{I assume that} {\textit{a}} {is not a preposition but a case marker that belongs to the DP (see \citealt{Lopez2000}, among others).}}{ quien tú quieras]]}

{  ‘you can invite everyone you want’}
\z

It is thus possible to analyse free relatives in Old Spanish as DPs under the assumption that empty syntactic categories correlate with certain semantic interpretations (see also \citealt{Rivero1988}):

\ea \label{ex:key:308} %ex 308 
{  [\textsubscript{Spec, TP} [\textsubscript{DP} \textsubscript{D°} ‘everyone’] [\textsubscript{CP1} qual\textsubscript{j} \textsubscript{TP1} pro \textsubscript{T°} quier [\textsubscript{CP2} t\textsubscript{j} que \textsubscript{TP2} pro \textsubscript{T°} lo viesse t\textsubscript{j}]]] \textsubscript{T°} hauria ende gran dolor].}

  ‘...who one wants that had seen it, would have had great pain and a lot of pity.’

{  (13th c.,}{ }{Anonymous,} {\textit{Gran Conquista}}{)}
\z

Once \textit{qual quier} lost its transparent relative clause structure, it was reanalysed as a determiner D° or a pronoun DP and filled the empty DP in an earlier stage (see \sectref{sec:5.3.1}):

\ea{} \label{ex:key:309} %ex 309 & ex 310
  [\textsubscript{TP} [\textsubscript{DP} [\textsubscript{CP} qual quier que lo viesse]] \textsubscript{T°} hauria ende gran dolor].

  [\textsubscript{TP} [\textsubscript{DP} qualquier (NP)] \textsubscript{T°} hauria ende gran dolor].

{  (my example for demonstration)}
\z

Let us now see how the syntactic change was interpreted semantically. A compositional analysis of Old Spanish example in \REF{ex:key:309} is presented in \REF{ex:key:311}. I assume that in Old Sp. \textit{qual} \textit{quier} means something like ‘which(ever) one wants (to choose)’. I underline the interpretation of the verb \textit{quier} and discuss it below in detail. The interpretation is derived step by step. First the relative clause is interpreted as an open proposition due to the wh-feature of \textit{qual} and the clause structure (CP) in \REF{ex:key:311-1}, then the relative clause is interpreted as a nominal modifier, that is, as a property of the noun N in \REF{ex:key:311-2}, then the DP is interpreted as universal quantification in \REF{ex:key:311-3}. The truth conditions of the sentence in which the DP occurs are spelled out in \REF{ex:key:311}.

\ea \label{ex:key:311} %ex 311
\ea{} \label{ex:key:311-1}   ⟦ [\textsubscript{CP} qualquier que lo viesse] ⟧ =

    would see it (x) ${\wedge}$ one wants to choose (x)\hfill\textit{open proposition}
\ex{} \label{ex:key:311-2}  ⟦ [\textsubscript{NP} Ø ‘person’ [\textsubscript{CP} qualquier que lo viesse]] ⟧ =

    λx[person (x) ${\wedge}$ would see it (x) ${\wedge}$ one wants to choose (x)] \hfill \textit{property}
\ex{} \label{ex:key:311-3}  ⟦ [\textsubscript{DP} [\textsubscript{NP} [\textsubscript{CP} qualquier que lo viesse]]] ⟧ =

    λP. \forall x[person (x) ${\wedge}$ one wants to choose (x) ${\wedge}$ would see it (x) ${\wedge}$ P(x)] \hfill  \textit{universal quantifier}
\ex{} \label{ex:key:311-4} ⟦qualquier que lo viesse hauria ende gran dolor⟧ = 1 if would have had′ (\forall x [person(x) ${\wedge}$ one wants to choose (x) ${\wedge}$ would see him (x)]) (great pain)

    Prose: For every person x that one wants to choose and that would see it, x would have great pain.
\z
\z

I will now discuss the semantics of \textit{qual quier} under the transparent meaning of the verb \textit{quiere} and the subject \textit{pro} at Stage 1 of the relative clause analysis. The most important question is whether this meaning is part of the truth conditions of the statement in \REF{ex:key:311}. I assume that it is not; that is, the open proposition \textit{x would have great pain} is true independently of the value one chooses for {x}. In other words, facts in the real world are independent of belief states. Thus, the statement \textit{Any even number (wichever one wants to choose) can be divided by 2} is true independent of whether someone wants to choose 4, 6, 8, etc. We can analyse the literal meaning of \textit{qual quier} rather as a side comment on the part of the speaker about the variable x, similar to appositive relative clauses as in \textit{My colleague (who is a nice guy) moved to Germany}. This side comment can be interpreted as a different illocution from the truth-conditional statement, namely as an invitation of the speaker to the addressee to choose any value for x. In this case, \textit{pro} is interpreted as a pronominal form of the addressee, as in \textit{pro} \textit{quisieres} ‘you (informal) would like’ in \REF{ex:key:312}. Note that in the example \REF{ex:key:312} \textit{qual quisieres} interrupts the sentence \textit{pon una estrella fixa por su longura}, which confirms the side comment analysis of \textit{qual quisieres}:

\ea \label{ex:key:312} %ex 312
  \gll damos te exiemplo a esto. pon una estrella fixa \textit{qual} \textit{quisieres} por su longura del eguador con el grado con que passa medio cielo.\\
give you example to this put a star fixed as wish for its length of.the equator with the degreee with which crosses halfway sky\\
  \glt ‘We are giving you an example of this. Put a fixed star (whichever you want) along the equator in an angle such that it passes through the middle of the sky.

{  (13th c., Alfonso X,} {\textit{Astronomía}}{)}
\z

The form \textit{qual quisier*} becomes less frequent in the 14th and 15th centuries and disappears completely in the 16th century. This is expected with the decrease of transparency of this variant of the construction.

In more formal texts, such as legal texts, where the addressee represents the general public or the reader of the law such as the context of \REF{ex:key:312}, the form \textit{qual quisieres} ‘which you (informal) want’ cannot be used. The subject \textit{pro} in \REF{ex:key:312} refers to the general public and the legal statement is true no matter which value the general public wants to choose for x, where x represents a person such that if x saw it, x would have great pain. According to the impersonal interpretation of \textit{pro}, it is guaranteed that the set of possible values for x is defined as large as possible (i.e. it contains all possible people the public can think of choosing for x). One direct consequence of this analysis of \textit{qual quier} as a side comment is that one should be able to find two verbs of \textit{querer}, one inside the proposition that has an effect on the truth conditions, and another inside a side comment that does not have any effect on the truth conditions. According to this analysis, only the embedded verb \textit{quisiere} in \REF{ex:key:313} has an effect on the truth conditions and \textit{qual quier} is interpreted as a side comment ‘whichever value one wants to choose for x’, where x is represented by \textit{qual}:

\ea \label{ex:key:313} %ex 313
  \gll Degrado deuen ser dadas las cosas espirituales \& no por preçio. \textit{onde} \textit{qual} \textit{quier} \textit{que} \textit{quisiere} \textit{entrar} \textit{en} \textit{orden} \textit{de} \textit{religion} \textit{no} \textit{deue} \textit{dar} \textit{precio} \textit{ninguno}\\
of.grace should be given the things spiritual \& \textsc{neg} for price where \textsc{qual} \textsc{quier} who whishes enter in order of religion \textsc{neg} should give price any\\
  \glt ‘Spiritual things must be given not bought. \textit{Whoever would want to enter in the realm of} \textit{religion should not give any price}.’

{  (13th c.,}{ }{Alfonso X,} {\textit{Siete partidas}}{)}
\z

The crucial assumption is that at the stage, when \textit{querer} was analysed as a lexical verb (i.e. before the 13th c. and partially also in 13th c.), the side comment expressed by \textit{qual quier} as ‘whichever one wants to choose’ or by \textit{qual quisieres} ‘whichever you, informal, would like to choose’ is literally expressed, whereas in later periods, when \textit{quier} is no longer interpreted as a verb, free choice is not literally encoded in the relative clause, but in the indefinite itself as part of its conventional meaning (see \chapref{sec:1} and example \REF{ex:key:4} repeated here in \REF{ex:key:314})

\ea \label{ex:key:314} %ex 314
  \gll {Puedes} {traer} \textit{cualquier} {libro.}\\
  can bring \textsc{cualquier} book\\
  \glt ‘You can bring any book.’\\
  Conventional meaning: ‘You can bring me a book, and each book is a possible option.’
  (\citealt{AloniEtAl2011}: 2)
\z

Another consequence of the change is that the choice to pick a value for x no longer depends on the generic subject \textit{pro} of \textit{quier} as in Old Spanish, but on the agent of the clause in which \textit{cualquiera} is used. In the Modern Spanish example in \REF{ex:key:314}, the subject of free choice is expressed by the grammatical subject \textit{pro}, which refers to the addressee, who can pick any book he/she wants, not the generic \textit{pro} inside \textit{cualquier}, which is no longer transparent in Modern Spanish. Another consequence of the grammaticalisation is that the indicative mood and perfective aspect increased (see \sectref{sec:5.2}). As was demonstrated in \sectref{sec:2.3}, \textit{cualquiera} has a ‘random’ interpretation or an evaluative interpretation under volitional verbs with perfective aspect (see \sectref{sec:2.3} and \chapref{sec:3}), repeated here in \REF{ex:key:315}. Under both interpretations (RCI and EVAL), the generic \textit{pro} is no longer visible as the agent of choice. In \REF{ex:key:315} the random choice depends on the explicit agent Juan and \REF{ex:key:315} does not express any choice at all, but speaker’s evaluation.

\ea \label{ex:key:315} %ex 315
  \gll Juan compró un libro \textit{cualquiera}.\\
  Juan bought a book \textsc{cualquiera}\\
  \glt ‘Juan bought a book.’
\ea  John bought a book, and he choose it indiscriminately. (Random Choice)
\ex John bought a book, and the book is not special or remarkable (according to the speaker). (EVAL)

    \citep[2]{am2018}
\z
\z

According to my analysis, the semantic change can be summarised in the following steps. The grammaticalisation of \textsc{qualquier} has induced a change in the interpretation of \textsc{qualquier} with respect to the subject of ‘want to choose a value for \textit{qual}’. When \textit{qualquier} was still transparent (before the 13th c. and partially still in 13th c.), free choice depended on the generic \textit{pro} and was expressed explicitly by the verbal phrase \textit{pro}\textsubscript{[generic]} \textit{quier}. After the loss of transparency, free choice was encoded in the indefinite and the agent of free choice depended on the grammatical subject/agent of the clause. Random Choice and EVAL appeared with the rise of indefinite nouns, indicative mood, and the volitional/predicative verb class:

\newpage
\ea \label{ex:key:316} %ex 316
  1.  \textit{qual} \textit{pro} \textsubscript{V°/T°} \textit{quiere que} V-subj. (before the 13th c.) $\rightarrow$

    Free Choice expressed explicitly and the agent of Free Choice is \textit{pro}\textsubscript{[generic]}

{} 2.{  Indefinite pronoun/determiner} {\textit{cualquier,-a}}{ (after the 13th c. till today) $\rightarrow$}

{    Free Choice expressed} {{implicitly}}{ as a conventional implicature and the agent of Free Choice is the agent of the matrix verb.}

  3.  \textsc{un} N \textit{cualquiera} under volitional verbs in indicative $\rightarrow$

{    Random Choice Modality (from the 18th c. until today)}

  4.  \textsc{un} N \textit{cualquiera} under predicative verbs in indicative (from the 18th c. till today)

EVAL
\z

Given the syntactic and semantic change suggested in \REF{ex:key:316}, the grammaticalisation of the indefinite does not correlate with a loss or weakening in meaning as \citet{cp2009} suggest — the grammaticalised form \textsc{qualquier} is not void of meaning — but it correlates with a change from an \textit{explicit} to an \textit{implicit} encoding of Free Choice, once \textit{qual quier} is no longer analysed literally (already in 13th c.); and it also correlates with the change of the agent of Free Choice from generic to grammatical subject/agent of the matrix clause. From the 13th century until today, Free Choice meaning remains, but it is encoded differently. Only much later, from the 18th century onwards, we observe that a {new meaning} (EVAL) appears due to the increase of indicative mood, predicative verbs, and the increase of \textsc{un} N \textit{cualquiera}. I will discuss the meaning change to EVAL in \sectref{sec:5.3.2}.

\subsection{{Analysis} {of} {\textsc{qualquier}} {as} {a} {postnominal} {modifier}} \label{sec:5.3.2}

The relative clause \textsc{qualquier} \textit{que…} could also appear postnominally with an overt noun as antecedent in the 13th century.

Let us now look into this type more closely. The antecedent of the relative clause introduced by \textit{qual(quier)} could not only appear with bare nouns, as shown in \sectref{sec:5.3.1} (which was the most frequent construction in the 13th century, as shown in \sectref{sec:5.2}), but also with overt quantificational determiners as in \REF{ex:key:317}, the frequency of which is very low in the 13th c.):

\ea \label{ex:key:317} %ex 317
\ea  \gll Mas si por \textit{alguna} \textit{otra} \textit{razon} \textit{qualquier} que no fuese delas sobre dichas enestas leyes deseredase el padre a su fijo […]\\
but if for some other reason \textsc{qualquier} that \textsc{neg} was of.the about mentioned in.these laws disinherits the father to his son\\
    \glt ‘But if for some other (whatsoever) reason, that is not among the ones mentioned before in these laws, the father disinherits his son […]’

    (13th c., Alfonso X, \textit{Siete Partidas})
    
\ex{} \gll [...] que ayan poder de tomar casas vinnas huertas oliuares molinos tiendas atahonas aldeas e \textit{todo} \textit{otro} \textit{heredamiento} \textit{qualquiere} que sea de todo omne [....].\\
{} that have power to take houses vineyards orchards olive.groves mills shops attics villages and all other inheritance \textsc{qualquiere} that is.\textsc{SUBJ} of all men\\
    \glt ‘[…] that they have the power to take houses, vineyards, vegetable gardens, olive groves, tents, horse mills, hamlets, and any other inheritance, whichever it might be, from every person […].’

    (13th c., Alfonso X, \textit{Doc. Cast. - Andalucía})
    
\ex \gll o por abenençia o por aluedrio o por conposicion o por \textit{toda} \textit{otra} \textit{manera} \textit{qualquiere} que el touiere por bien\\
or by benevolence or by free.will or by composition or by any other way \textsc{qualquiere} that he considers for good\\
\glt ‘by the agreement or by arbitral award or by a deal or by any other manner that he approves of’

(13th c., Alfonso X, \textit{Doc. Cast. - Castilla la Nueva})
    
\ex{} \gll […] Desi dixo domjngujllo que le diese el Rey \textit{vn} \textit{omne} \textit{qual} \textit{quier} a que le firiese e quando le oujese ferido que se yrie para el castillo\\
{} thus said Domjngujllo that him give the King a man \textsc{qual} \textsc{quier} to whom him would.strike and when him been.heard wounded that \textsc{refl} should.go to the castle\\
    \glt ‘Then Dominguillo said thet the King should give him a man which(ever) one wants whom he could hurt, and after hurting him, he would go to the castle.’

    (14th c., Anonymous, \textit{Veinte Reyes})
\z
\z

This shows that \textit{qualquier} cannot be interpreted as a quantificational determiner at this stage yet. Otherwise we would end up with two different quantifiers that quantify over the same NP or over the same domain. This is impossible in natural language; consider, for instance, English \textit{*some whichever reason} or \textit{*every whichever inheritance} or \textit{*every whichever manner} (see also \chapref{sec:2} on the exclusion of double quantification). Let us analyse the example of \textit{todo otro heredamiento que} in \REF{ex:key:317}, of which I only present the relevant part in \REF{ex:key:318}:

\ea \label{ex:key:318} %ex 318
  \gll ...que ayan poder de tomar casas, ... e \textit{todo} \textit{otro} \textit{heredamiento} \textit{qualquiere} \textit{que} \textit{sea}\\
that may.have power of take houses, {} and all other inheritance \textsc{qualquiere} that may.be\\
  \glt ‘…that they can take every inheritance that one wants it to be’
\z

The interpretation is derived step by step in \REF{ex:key:319}. First, the relative clause is interpreted as an open proposition due to the wh-feature of \textit{qual} and the clause structure (CP), then the relative clause is interpreted as a nominal modifier, that is, as a property, then the DP is interpreted as universal quantification. The truth conditions of the sentence in which the DP occurs are spelled out in \REF{ex:key:319-4}.\footnote{{Another
      possibility would be to interpret the postnominal} {\textsc{qualquier}} {as a free relative clause without an overt antecedent. On this analysis, one could interpret the free relative clause as a proposition with a discourse variable and not as a property that modifies an antecedent NP:}

      \ea
        {Some other thing. Whatever it [=this thing] might be.}
      \z
      {I leave it for future research, to distinguish between the two options.}

}

\ea \label{ex:key:319}
\ea{} \label{ex:key:319-1}  ⟦qualquier que sea⟧ =

    (i) there might be(x) \hfill open proposition
\ex{} \label{ex:key:319-2} ⟦heredamiento qualquier que sea⟧ =

    (ii) λx[inheritance (x) ${\wedge}$ there might be (x)] \hfill property
\ex{} \label{ex:key:319-3} ⟦todo otro heredamiento qualquier que sea⟧ =

    (iii) λP. All x[other (x) ${\wedge}$ inheritance(x) ${\wedge}$ there might be ${\wedge}$ P(x)]  universal quantifier
\ex{} \label{ex:key:319-4} ⟦\REF{ex:key:318} = que ayan poder de tomar todo otro heredamiento qualquier que sea⟧ = 1 iff

    can take (\forall x [inheritance(x) ${\wedge}$ there might be (x)]) (they)

    prose: according to the law they can take every inheritance that there might be
\z
\z

The only difference between headed relatives as in \REF{ex:key:319} and free relatives suggested in \sectref{sec:5.3.2} is that the nominal antecedent is overtly present in headed relatives in contrast to free relatives with a covert nominal antecedent in \sectref{sec:5.3.2}. The function of the relative clause CP \textit{qual quier que sea} is the same. It modifies an (overt) nominal antecedent and is thus of a property type <e,t>, similar to adjectives or other nominal modifiers:

\ea \label{ex:key:320} %ex 320
  \textit{Qualquier} with an overt antecedent:

  a. [\textsubscript{DP}\textit{\textsubscript{<{}<e,t>,t>}} alguna/toda [\textsubscript{Spec,NP} otra \textsubscript{<e,t>} [\textsubscript{NP <e,t>} razon/manera [\textsubscript{CP<e,t>} [\textsubscript{Spec,CP} [\textsubscript{D°} \textit{qualquier}\textsubscript{j}] \textsubscript{C′ C°} \textit{que} \textsubscript{TP2} \textit{pro\textsubscript{} }\textsubscript{T°} \textit{sea} t\textsubscript{j}]]]]

%%[Warning: Draw object ignored]
\z

I will show now what happened after \textit{qual quier que} was no longer analysed as a relative clause in headed relatives. In a nutshell, the idea is that the string \textit{qual quier} lost its transparent status as consisting of the wh-pronoun \textit{qual} and the verbal form \textit{quier} (see \sectref{sec:5.3.1}); that is, the wh-feature disappeared, and the verb was no longer analysed as such, with the consequence that the embedded clause introduced by \textit{que} was no longer necessary. This is why the frequency of \textit{qualquier,-e,-a que}… decreases, as already shown in \sectref{sec:5.2}. The postnominal \textit{qualquier,-e,-a} thus can appear without the embedded clause (see, e.g. \ref{ex:key:321}):

\ea \label{ex:key:321} %ex 321
    \gll o con cuchillo o con \textit{otra} \textit{cosa} \textit{qualquier}\\
or with knife or with other thing \textsc{qualquier}\\
    \glt ‘or with a knife or with some other thing whichever one wants’

    (13th c., Alfonso X, \textit{Siete Partidas})
\z

The crucial question is what syntactic category and semantic interpretation the postnominal \textit{qualquier,-e,-a} acquired in structures like \REF{ex:key:321}, once it was no longer analysed transparently.

I assume that the postnominal \textit{cualquiera} has preserved its syntactic status as a nominal modifier (like relative clauses) and its semantic status as a property of type <e,t>, as shown in \REF{ex:key:322}. The only difference to relative clauses like those in \REF{ex:key:320} is that the new postnominal \textsc{cualquier} is no longer clausal, but a complex compound word (for details, see \chapref{sec:3}):

\ea{} \label{ex:key:322}
  [\textsubscript{DP} una [\textsubscript{NP} cosa/razon [\textsubscript{ModifP} \textit{cualquiera}]]]

%%[Warning: Draw object ignored]
\z

As I already pointed out, postnominal \textit{cualquiera} does not have the evaluative meaning before the 17th century (see \sectref{sec:5.1} and \sectref{sec:5.2}). The free choice interpretation of the postnominal modifier is triggered by some modal verb, subjunctive, or generic operator (see \chapref{sec:2}):

\newpage
\ea \label{ex:key:323}
  postnominal \textsc{cualquiera}

  \gll ...sea debajo de la equinoccial o (sea) en \textit{otra} \textit{parte} \textit{cualquiera}» A este propósito, trae Celio Rodiginio lo de Arriano, historiador griego\\
...be under of the equinoctial or (be) in other part \textsc{cualquiera} to this purpose, brings Celio Rodiginio the of Arriano, historian Greek\\
  \glt ‘… be it under the equinoctial or (be it) in any other place» In this respect, Celio Rodiginio says of Arriano, Greek historian’ (brackets my own)

  (16th c., Torquemada, \textit{Jardín})
\z

The free choice meaning of the postnominal modifier \textit{cualquiera} in \REF{ex:key:324} is similar to the meaning of a postnominal relative clause \textit{qual quiera} at earlier stages (see \sectref{sec:5.3.1}). To illustrate this example, here it can be roughly paraphrased as \textit{some place x}, where x can have any value (one wants). The semantic shift from FCI to EVAL in the 17th century in \REF{ex:key:324} can be analysed as a shift from free choice \textit{cualquiera} as ‘any (one wants)’ to a property \textit{cualquiera} that ‘anyone has’ and thus to a property with the meaning ‘ordinary’, because this property does not distinguish individuals from each other (see \chapref{sec:3} for details):

\ea \label{ex:key:324} %ex 324
  \gll Juan es un hombre \textit{cualquiera} (para mí)\\
Juan is a man \textsc{cualquiera} (for me)\\
  \glt ‘Juan is a man with a property that anyone has’ (literal) $\rightarrow$

‘Juan is an ordinary man (to me).’

  (@ 151)
\z

The evaluative interpretation has triggered a further shift of postnominal modifying \textit{cualquiera} into nominal \textit{cualquiera} as in \textit{es una cualquiera}, which appears late in the CdE from the 18th century (from 14th to 18th\textsuperscript{} c. 0.0 per mil in each century, 18th\textsuperscript{} c. 0.20 per mil, 19th\textsuperscript{} c. 0.73 per mil). Note that, semantically, this nominal \textit{cualquiera} is still interpreted as an evaluative predicate of the type <e,t>, but behaves differently from the postnominal function in that it also encodes the feature [+animate] and female gender:

\ea \label{ex:key:325} %ex 325
 \gll \textit{¡ya} \textit{se} \textit{ve!}, \textit{es} \textit{una} \textit{cualquiera};\\
  already \textsc{refl} sees is a \textsc{cualquiera}\\
\glt ‘One can already see! She is a promiscuous woman;’

  (19th c., Rizal, \textit{Noli me tangere})
\z

Recall that, in Argentinian Spanish, the semantic shift into the evaluative predicate has triggered a further syntactic reanalysis of \textsc{cualquiera} a gradable adjective (see \chapref{sec:4}):

\ea \label{ex:key:326} %ex 326
  \gll Juan es muy \textit{cualquiera} (para mí)  \hfill (ArgSp)\\
  Juan is very \textsc{cualquiera} for me\\
  
  {[ordinary′ (Juan) (to me) (to a high degree)]}\\
  \glt ‘Juan is very common/ordinary to me.’

    (@ 152)
\z

I assume that it is the semantic interpretation of relative clauses as properties that enables one to interpret \textit{cualquiera} as a property with an evaluative meaning ‘ordinary’ in predicative position such as \textit{es un hombre cualquiera} ‘he is an ordinary man’ or as an evaluative property of women that has been lexicalised in \textit{una cualquiera} (\chapref{sec:3} on EVAL). The rise of the indicative mood and predicative verbs has further enhanced the property interpretation of \textit{cualquiera} as in N \textit{cualquiera} (see \chapref{sec:3}).

The different analyses of prenominal and postnominal \textit{cualquiera} explains why only the postnominal \textit{cualquiera} has different lexical meanings (FCI, EVAL) in European Spanish, but not prenominal \textit{cualquiera}. The reason for this is that only postnominal \textit{cualquiera} has a predicative status and is interpreted as a nominal modifier similar to adjectives, whereas the quantificational determiner is a functional category that is more restricted in meanings than a lexical category.

\section{{Summary} {and} {discussion}} \label{sec:5.4}

I have proposed an analysis that shows how the reanalysis of \textsc{qualquier} as a relative clause into a quantificational determiner (which \citealt{cp2009} and others call grammaticalisation of \textsc{qualquier}) has taken place, namely from a relative clause to reanalysis as a word compound \textsc{qualquier} (without the ending \textit{-a}), which was categorised as a quantificational determiner D° with a FC interpretation ‘any(one/thing)’. I have provided a detailed analysis of the reanalysis process that involved a loss of compositional meaning of the generic subject pronoun, the modal verb \textit{quier} and the relative wh-pronoun \textit{qual} and the appearance of a single word with a free choice interpretation associated with the matrix subject. The steps and empirical evidence is summarised as follows.

%In this section, I have shown on the basis of numerically quantified and statistically analysed data how \textsc{qualquier} has changed on the syntactic and semantic level. I have suggested an analysis that shows how the reanalysis of \textsc{qualquier} as a relative clause into a quantificational determiner (which \citealt{cp2009} and others call grammaticalisation of \textsc{qualquier}) has taken place, namely from a relative clause to lexicalisation as the single word \textsc{qualquier} (without the ending \textit{{}-a}), which was subsequently reanalysed as a quantificational determiner D° with FC interpretation ‘any(one/thing)’:

\ea \label{ex:key:327} %ex 327
  Free relative clause \textit{qual quier} was reanalysed as an indefinite determiner (grammaticalisation as a determiner, see also the Innovation/Grammaticalisation Hypothesis in §2)
  \begin{itemize}
  \item Free Rel. Clause $\rightarrow$ Determiner as in \textit{cualquier} N
  \item Evidence:
  \begin{itemize}
    \item The separated forms in which a noun or an NP can be placed between relative pronoun \textit{qual} and verb \textit{quier} (e.g. \textit{qual manera quier} lit. ‘which manner one wants’) from 13th c.
    \item After 13th c. \textit{qual quier} N(P) ‘which one wants N(P)’ continued the path of grammaticalisation, becoming more fixed, as shown by graphically synthetic form \textit{qualquier} and later \textit{cualquier}.
    \item The rise in frequency of \textit{qualquier} NP, the decrease in use of relative clause markers such as \textit{que} ‘that’ is another support for the grammaticalisation.
  \end{itemize}
  \item Interpretation of Determiner \textit{cualquier}: FCI only
  \end{itemize}

\z

I have suggested that the change of postnominal \textsc{qualquier} (with the ending \textit{{}-a}) was a change from a headed relative clause into a postnominal modifier as in \textsc{UN} N \textit{cualquiera} with the subsequent change into a noun \textit{cualquiera} as in \textsc{UN} \textit{cualquiera} in European Spanish. The steps and evidence are summarised as follows.

\ea \label{ex:key:328}
  Headed Rel. Clause [N \textit{qual quiera}] was reanalysed as a nominal modifier \textit{qualquiera}
    \begin{itemize}
    \item NP Rel. Clause $\rightarrow$ bare N \textit{qualquier} $\rightarrow$ bare N \textit{qualquier} $\rightarrow$ Article N \textit{cualquiera} $\rightarrow$ Article \textit{cualquiera}
    \item Evidence: 
      \begin{itemize}
      \item graphically synthetic form \textit{qualquier,e,a} and later \textit{cualquiera},
      \item the increase of Art N \textit{cualquiera}
      \item the increase of ambiguity between nominal modifiers and adjectives,
      \item the increase in lexical meanings (FCI, EVAL),
      \item the increase in UN \textit{cualquiera} (EVAL)
      \item the increase in definite Art N \textit{cualquiera} (EVAL)
      \end{itemize}
    \end{itemize}
\z

Recall from \chapref{sec:2} that \citet{cp2009} assume that \textsc{qualquier} has gradually lost its semantic interpretation under the assumption that grammaticalisation correlates with a loss or weakening of semantic content (\citealt{Lehmann2002}, \citet{cp2009}). At the same time,  \citet{cp2009} assume that \textsc{qualquier} undergoes a “subjectification” process going from “election” to “generalisation” to “valorisation”. These two processes (grammaticalisation and subjectification) do not represent a contradiction, if we take the syntactic nature of \textsc{qualquier} into account and the linguistic context in which it appears (e.g., indicative and modal verbal mood). It’s the postnominal \textsc{qualquier} under indicative verbal mood, not the prenominal \textsc{qualquier}, that undergoes the "subjectification" process and acquires the new meaning EVAL in European Spanish, which triggers a further process of syntactic reanalysis into a gradable adjective as in \textit {re/medio cualquiera} in Argentinian Spanish (\chapref{sec:4}). The postnominal \textsc{qualquier} certainly has grammaticalised to some degree, as it lost the transparent relative clause interpretation of the wh-word qual + modal verb \textit{querer} and was reanalysed as a single word (see \cite{Lehmann2002} on prototypical features of grammaticalisation). However, the grammaticalisation of the postnominal \textsc{qualquier} into a nominal modifier that is similar to the syntactic category of adjectives is what triggered further linguistic changes such as the appearance of articles and meaning ambiguity depending on the linguistic context (+/- indicative verbal mood). The two elements -- the determiner \textsc{qualquier} and the postnominal modifier \textsc{qualquier} -- followed a different linguistic path. One could also describe the path of the postnominal modifier \textsc{qualquier} in the diachrony of Spanish as “constructionalisation” (\citealt{traugott2013constructionalization}) because it has acquired both new form and new meaning, which is particularly seen in Argentinian Spanish where it has acquired the new form of a gradable adjective with the meaning ‘ordinary’ and the new form of an evaluative noun with the meaning ‘bullshit/non-sense’.

%If \citealt{cp2009} mean that \textsc{qualquier} lost the literal meaning of \textit{querer} ‘want’, I agree with this observation. The crucial change according to my analysis is the change in the agent of Free Choice Modality (i.e. from the generic \textit{pro} of verb \textit{quier} to the agent of the matrix verb). However, after the recategorization of the relative clause as determiner in the case of prenomoninal \textsc{qualquier} and nominal modifier in the case of postnominal \textsc{qualquier}, further semantic and syntactic changes appeared in European and in Argentinian Spanish. The postnominal modifier received a kind property visible with definite articles as in \textit {el hombre cualquiera} ‘the common man’ in the modern use of European Spanish and it was reanalysed as a degree adjective with the meaning ‘unremarkable, bad, etc.’ in Argentinian Spanish.

%Unlike what is usually claimed for grammaticalisation, indefinites such as \textsc{qualquier} do not lose semantic content. The reason may be that this is not a prototypical case of grammaticalisation. \textsc{qualquier} certainly has grammaticalised to some degree, as it lost the transparent relative clause interpretation of the wh-word \textit{qual} + modal verb \textit{querer} and was reanalysed as a single word (see \citealt{Lehmann2002} [1982] on prototypical features of grammaticalisation). However, further semantic and syntactic changes of (postnominal) \textit{cualquiera} in European and of \textit{cualquiera} in Argentinian Spanish appeared. Strictly speaking, the postnominal \textit{cualquiera} did not grammaticalise into a functional category like the prenominal \textit{qual quier} did in European Spanish in the 13th century. This is probably why postnominal \textit{cualquiera} and prenominal \textit{cualquier} differ in syntax and semantic interpretations in European Spanish.

Finally, I would like to discuss the concepts of grammaticalisation and lexicalisation in order to understand which of these processes are involved in the linguistic changes of \textsc{qualquier}. 
Many studies have recognised the challenge of distinguishing the processes of grammaticalisation and lexicalisation (e.g. \citealt{BrintonTraugott2005,diewald2012paradigmatic,Lehmann2002,traugott2013constructionalization,trousdale2012grammaticalization,van2015lexicalization}). Compound words like \textit{going to, a lot of, cualquiera} are among the most challenging cases with regard to the question of which of the two processes (lexicalisation or grammaticalisation) is involved in the linguistic change of these words (e.g. \citealt{BrintonTraugott2005,trousdale2012grammaticalization,van2015lexicalization}).  Both processes can involve similar linguistic mechanisms such as reanalysis, desemanticisation, and change in syntactic categorisation. Regarding reanalysis, linguistic elements can undergo reinterpretation in the process of grammaticalisation and lexicalisation. Grammaticalisation often involves a reduction in semantic content as the word compound qualquier transitions into a grammatical category (quantificational indefinite). Lexicalisation can also involve semantic bleaching, where a compound phrase becomes a single lexical item. Both processes can involve changes in the grammatical category. In grammaticalisation, words often lose their full lexical category status, while in lexicalisation, phrases can solidify into a single lexical category. This overlap suggests that grammaticalisation and lexicalisation are not completely distinct processes. The term lexicalisation has at least two meanings. It means the incorporation of a complex linguistic expression such as relative clause \textit {qual quier} into the lexicon as a word(compound) (meaning 1) and it also means a reanalysis of a complex expression as a lexical category with lexical content (e.g. the complex expression \textit {valer la pena} in Spanish means ‘to be worth’) in contrast to a grammatical or functional category with a grammatical function (e.g. determiners) (meaning 2). This said, the process of grammaticalisation of \textit {qual quier} necessarily involves lexicalisation in the first sense of the term (meaning 1). The second sense of the term lexicalisation only applies to the postnominal modifier \textit{cualquiera} that has been reanalysed into a lexical category nominal modifier or a noun with lexical meanings such as ‘ordinary’ as in \textit{un día cualquiera} or ‘male nobody’ as in \textit{un cualquiera}.      

